% main.tex -- "Architecting Computational Mechanics".
% Build: latexmk -pdf main     (or pdflatex main && bibtex main && pdflatex main && pdflatex main)
% Numbers in the text are macros from tables/*.tex, generated by scripts/*.py.
\documentclass[11pt,a4paper,twoside,openany]{book}
% macros.tex -- preamble of the book "Architecting Computational Mechanics" (A4, class book).
\usepackage[T1]{fontenc}
\usepackage[utf8]{inputenc}
\usepackage{lmodern}
\usepackage[a4paper,margin=2.6cm]{geometry}
\usepackage{amsmath,amssymb,bm,amsthm}
\usepackage{graphicx}
\usepackage{booktabs,array,longtable,tabularx}
\usepackage{xcolor}
\usepackage{microtype}
\usepackage{enumitem}
\usepackage{xspace}
\usepackage{subcaption}
\usepackage[most]{tcolorbox}
\usepackage{listings}
\usepackage{tikz}
\usetikzlibrary{positioning,fit,arrows.meta,calc,backgrounds}
\usepackage[sectionbib,numbers,sort&compress]{natbib}
\usepackage{chapterbib}
\renewcommand{\cite}{\citep}
\setlength{\bibsep}{2pt}
\renewcommand{\bibsection}{\section*{References}\addcontentsline{toc}{section}{References}}
% every chapter ends with \chapterrefs: its own list of references, numbered from [1]
\newcommand{\chapterrefs}{\bibliographystyle{unsrtnat}\bibliography{refs}}
\usepackage{url}\urlstyle{tt}
\makeatletter\def\UrlBreaks{\do\_\do\.\do\/\do\,}\def\UrlBigBreaks{\do\:\do\@}\makeatother
\usepackage[colorlinks=true,linkcolor=blue!50!black,citecolor=blue!50!black,urlcolor=blue!50!black]{hyperref}
\usepackage[capitalise,noabbrev]{cleveref}
\setlist{itemsep=2pt,topsep=3pt,parsep=0pt}
\setlength{\emergencystretch}{5em}
\usepackage{fancyhdr}\pagestyle{fancy}\fancyhf{}\fancyhead[LE,RO]{\thepage}\fancyhead[RE]{\small\nouppercase{\leftmark}}\fancyhead[LO]{\small\nouppercase{\rightmark}}\renewcommand{\headrulewidth}{0.3pt}\fancypagestyle{plain}{\fancyhf{}\fancyfoot[C]{\thepage}\renewcommand{\headrulewidth}{0pt}}
\numberwithin{equation}{chapter}
\numberwithin{figure}{chapter}
\numberwithin{table}{chapter}

% names and code references
\newcommand{\fea}{\texttt{fea\_engine}}
\newcommand{\rom}{\texttt{rom\_engine}}
\DeclareRobustCommand{\code}[1]{\path{#1}}
% code anchor: [Code: importable.name @ tag]; the tag is set once for the whole book
\newcommand{\booktag}{v1.0.1\,/\,v0.2.0}
\DeclareRobustCommand{\anchor}[1]{{\sffamily[Code:\,}\path{#1}{\sffamily]}}

% theorem-like environments
\theoremstyle{plain}
\newtheorem{proposition}{Proposition}[chapter]
\theoremstyle{definition}
\newtheorem{definition}{Definition}[chapter]
\newtheorem{example}{Example}[chapter]
\theoremstyle{remark}
\newtheorem{remark}{Remark}[chapter]

% boxes
\tcbset{boxsep=3pt,left=5pt,right=5pt,top=3pt,bottom=3pt,fonttitle=\bfseries\small,fontupper=\small,breakable}
\newtcolorbox[auto counter,number within=chapter]{algbox}[2][]{colback=gray!4,colframe=gray!55,title={Algorithm \thetcbcounter: #2},#1}
\newtcolorbox{limitsbox}[1][Limits, as measured]{colback=orange!5,colframe=orange!60!black,title={#1}}
\newtcolorbox{codebox}[1]{colback=blue!3,colframe=blue!40!black,title={#1}}

% listings
\lstdefinestyle{py}{language=Python,basicstyle=\ttfamily\footnotesize,keywordstyle=\color{blue!60!black}\bfseries,
  commentstyle=\color{green!40!black}\itshape,stringstyle=\color{red!60!black},showstringspaces=false,
  frame=single,framerule=0.3pt,rulecolor=\color{gray!60},numbers=left,numberstyle=\tiny\color{gray},numbersep=5pt,
  breaklines=true,columns=fullflexible,keepspaces=true,xleftmargin=1.5em,framexleftmargin=1.2em,
  belowskip=0.6em,aboveskip=0.6em}
\lstset{style=py}

% evidence: generated numbers (scripts/), and the status of a chapter
\newcommand{\evidence}[1]{\textit{Evidence: #1}}
\newcommand{\machine}{the sandbox used for the figures of this chapter (2 CPUs)}
\newcommand{\todo}[1]{\textcolor{red!70!black}{\textbf{[to do: #1]}}}
\newcommand{\stubchapter}[1]{%
\begin{tcolorbox}[colback=yellow!8,colframe=yellow!50!black,title={Not yet written}]
This chapter is planned in \texttt{plan/main.pdf} (#1). It will follow the chapter template of the plan and the structure of Chapter~\ref{ch:contact}.
\end{tcolorbox}}

\newcommand{\KKT}{\mathrm{KKT}}
\DeclareMathOperator{\argmin}{arg\,min}
% generated numbers (scripts/*.py); do not edit

% generated by scripts/ch01_assets.py -- do not edit
\newcommand{\cdeRegTotal}{34\xspace}
\newcommand{\cdeRegPass}{34\xspace}
\newcommand{\cdeRegBad}{none\xspace}
\newcommand{\cdeFeaImportsRom}{0\xspace}
\newcommand{\cdeRomImportsFea}{0\xspace}
\newcommand{\cdeRomTestsImportFea}{35\xspace}
\newcommand{\cdeRomTestFiles}{85\xspace}
\newcommand{\cdeExtStatic}{$3.1\times10^{-15}$\xspace}
\newcommand{\cdeExtFreqErr}{$2.0\times10^{-4}$\xspace}
\newcommand{\cdeExtFreqOrder}{2.0\xspace}
\newcommand{\cdeExtFreqCoarse}{27\%\xspace}
\newcommand{\cdeDispSlope}{2.3\xspace}
\newcommand{\cdeDenseFiveK}{0.2\xspace}
\newcommand{\cdeDenseFiftyK}{20.0\xspace}
\newcommand{\cdeDispMaxDof}{3\,362\xspace}
\newcommand{\cdeDispThreshold}{5\,000\xspace}
\newcommand{\cdeDispExtraMs}{531.4\xspace}
\newcommand{\cdeDispExtraFrac}{1790\xspace}
\newcommand{\cdeDispExtraFracSmall}{56\xspace}
\newcommand{\cdeDispExtraFracLarge}{1790\xspace}
\newcommand{\cdeHasTorch}{not installed\xspace}
\newcommand{\cdeMachine}{Linux 6.8.0-138-generic, 2 CPUs, NumPy 2.2.6\xspace}

% generated by scripts/ch02_assets.py -- do not edit
\newcommand{\cmeGenAllPositive}{yes\xspace}
\newcommand{\cmeGenNRows}{7\xspace}
\newcommand{\cmeRectErr}{0\xspace}
\newcommand{\cmeHoleErr}{$5.7\times10^{-3}$\xspace}
\newcommand{\cmeHoleOffErr}{$7.9\times10^{-5}$\xspace}
\newcommand{\cmeNotchErr}{0\xspace}
\newcommand{\cmeFilletErr}{$2.3\times10^{-4}$\xspace}
\newcommand{\cmeBoxHoleErr}{$1.2\times10^{-2}$\xspace}
\newcommand{\cmeArcOrder}{2.2\xspace}
\newcommand{\cmeArcErrCoarse}{$2.2\times10^{-2}$\xspace}
\newcommand{\cmeArcErrFine}{$1.8\times10^{-4}$\xspace}
\newcommand{\cmeGradedSizeRatio}{4.6\xspace}
\newcommand{\cmeGradeForTarget}{1.14\xspace}
\newcommand{\cmeGradeRealTarget}{1.199\xspace}
\newcommand{\cmeGradeWorstOver}{-0.01\xspace}
\newcommand{\cmeGradeRatioPTwo}{3.00\xspace}
\newcommand{\cmeGradeRatioPOne}{1.00\xspace}
\newcommand{\cmeMeshioAllSet}{yes\xspace}
\newcommand{\cmeMBQuads}{16\xspace}
\newcommand{\cmeMBTris}{32\xspace}
\newcommand{\cmeMBShared}{5\xspace}
\newcommand{\cmeMBErr}{$6.8\times10^{-16}$\xspace}
\newcommand{\cmeMachine}{Linux 6.8.0-138-generic, 2 CPUs, NumPy 2.2.6\xspace}

% generated by scripts/ch03_assets.py -- do not edit
\newcommand{\cecPatchWorst}{$1.7\times10^{-14}$\xspace}
\newcommand{\cecPatchCount}{8\xspace}
\newcommand{\cecBendQuadFour}{0.667\xspace}
\newcommand{\cecBendQuadThirtytwo}{0.992\xspace}
\newcommand{\cecBendTriFour}{0.223\xspace}
\newcommand{\cecBendTriThirtytwo}{0.948\xspace}
\newcommand{\cecBendQuadEightErr}{$8.0\times10^{-13}$\xspace}
\newcommand{\cecBendTriSixErr}{$5.6\times10^{-14}$\xspace}
\newcommand{\cecBendQuadOrder}{1.9\xspace}
\newcommand{\cecLockHexPlainA}{0.71\xspace}
\newcommand{\cecLockHexPlainB}{0.18\xspace}
\newcommand{\cecLockHexBbarA}{0.73\xspace}
\newcommand{\cecLockHexBbarB}{0.74\xspace}
\newcommand{\cecLockHexTwentyA}{0.99\xspace}
\newcommand{\cecLockHexTwentyB}{0.90\xspace}
\newcommand{\cecEBTipErr}{$6.7\times10^{-16}$\xspace}
\newcommand{\cecEBFreqOrder}{4.0\xspace}
\newcommand{\cecTimoOrder}{2.0\xspace}
\newcommand{\cecEBFreqHz}{20.95\xspace}
\newcommand{\cecTimoShearShare}{0.65\xspace}
\newcommand{\cecMassConsSign}{positive\xspace}
\newcommand{\cecMassLumpSign}{negative\xspace}
\newcommand{\cecMassConsFine}{$8.0\times10^{-4}$\xspace}
\newcommand{\cecMassLumpFine}{$2.4\times10^{-3}$\xspace}
\newcommand{\cecMassConsOrder}{2.0\xspace}
\newcommand{\cecMassLumpOrder}{2.0\xspace}
\newcommand{\cecPlateSriFine}{1.078\xspace}
\newcommand{\cecPlateFullCoarse}{0.005\xspace}
\newcommand{\cecPlateFullFine}{0.602\xspace}
\newcommand{\cecPlateShellFine}{1.078\xspace}
\newcommand{\cecPlateSriCoarse}{1.059\xspace}
\newcommand{\cecPlateShellCoarse}{1.051\xspace}
\newcommand{\cecPlateRatioFS}{0.9997\xspace}
\newcommand{\cecMachine}{Linux 6.8.0-138-generic, 2 CPUs, NumPy 2.2.6\xspace}

% generated by scripts/ch04_assets.py -- do not edit
\newcommand{\cdyNewmarkOrder}{1.98\xspace}
\newcommand{\cdyNewmarkErrCoarse}{$5.8\times10^{-1}$\xspace}
\newcommand{\cdyNewmarkErrFine}{$2.4\times10^{-3}$\xspace}
\newcommand{\cdyNewmarkAmpWorst}{5\%\xspace}
\newcommand{\cdyFirstFreq}{20.951\xspace}
\newcommand{\cdyDtcLumped}{$5.0\times10^{-5}$\xspace}
\newcommand{\cdyDtcConsistent}{$2.7\times10^{-5}$\xspace}
\newcommand{\cdyDtcRatio}{1.83\xspace}
\newcommand{\cdyDtcLocal}{1.00\xspace}
\newcommand{\cdyGrowStable}{1.01\xspace}
\newcommand{\cdyGrowUnstable}{$3.1\times10^{154}$\xspace}
\newcommand{\cdyGrowJust}{0.99\xspace}
\newcommand{\cdyGrowJustOver}{$1.6\times10^{49}$\xspace}
\newcommand{\cdyPeakAll}{1.932\xspace}
\newcommand{\cdyPeakOne}{1.911\xspace}
\newcommand{\cdyPeakSDOF}{1.969\xspace}
\newcommand{\cdyPeakRatioOne}{0.971\xspace}
\newcommand{\cdyModalErrOne}{$2.9\times10^{-2}$\xspace}
\newcommand{\cdyModalErrFour}{$5.2\times10^{-4}$\xspace}
\newcommand{\cdyModalErrEight}{$7.7\times10^{-5}$\xspace}
\newcommand{\cdyModalAll}{39\xspace}
\newcommand{\cdyPeakNewmark}{1.932\xspace}
\newcommand{\cdyRayleighMid}{0.49\xspace}
\newcommand{\cdyHarmZeroErr}{$6.2\times10^{-13}$\xspace}
\newcommand{\cdyShareOne}{0.971\xspace}
\newcommand{\cdyHarmPred}{48.5\xspace}
\newcommand{\cdyHarmDiff}{$7.6\times10^{-8}$\xspace}
\newcommand{\cdyHarmPeakHz}{20.9\xspace}
\newcommand{\cdyHarmAmp}{49\xspace}
\newcommand{\cdyQuality}{50\xspace}
\newcommand{\cdyRandomDiff}{$3.2\times10^{-9}$\xspace}
\newcommand{\cdyMachine}{Linux 6.8.0-138-generic, 2 CPUs, NumPy 2.2.6\xspace}

% generated by scripts/ch05_assets.py -- do not edit
\newcommand{\cnlObjTLMax}{$4.2\times10^{-13}$\xspace}
\newcommand{\cnlObjLinTen}{30.69\xspace}
\newcommand{\cnlObjLinNinety}{2020.31\xspace}
\newcommand{\cnlObjLinMax}{4040.61\xspace}
\newcommand{\cnlTangentWorst}{$9.1\times10^{-11}$\xspace}
\newcommand{\cnlTrussPeak}{15926\xspace}
\newcommand{\cnlTrussPeakDelta}{42.26\xspace}
\newcommand{\cnlLoadCtrlFailed}{fails\xspace}
\newcommand{\cnlTrussErrDC}{$6.8\times10^{-14}$\xspace}
\newcommand{\cnlTrussErrAL}{$1.1\times10^{-13}$\xspace}
\newcommand{\cnlTrussALpts}{20\xspace}
\newcommand{\cnlPlPeak}{1.80\xspace}
\newcommand{\cnlPlReverseYield}{1.80\xspace}
\newcommand{\cnlPlHardenedYield}{1.8\xspace}
\newcommand{\cnlPlErr}{$3.2\times10^{-15}$\xspace}
\newcommand{\cnlPlUnloadSlope}{1.000\xspace}
\newcommand{\cnlPlTangentSlope}{4.88\xspace}
\newcommand{\cnlPlMin}{-3.03\xspace}
\newcommand{\cnlLCdelta}{22.302\xspace}
\newcommand{\cnlLCerr}{$1.1\times10^{-13}$\xspace}
\newcommand{\cnlFesPeakErr}{$1.6\times10^{-16}$\xspace}
\newcommand{\cnlFesSetErr}{$1.7\times10^{-16}$\xspace}
\newcommand{\cnlFesSet}{40.000\xspace}
\newcommand{\cnlMachine}{Linux 6.8.0-138-generic, 2 CPUs, NumPy 2.2.6\xspace}

% generated by scripts/ch06_assets.py -- do not edit
\newcommand{\cbkChunkDiff}{$2.0\times10^{-16}$\xspace}
\newcommand{\cbkModes}{q: batched\xspace}
\newcommand{\cbkReasonPlastic}{TrussPlastic2D does not declare supports\_batch\xspace}
\newcommand{\cbkChunkedPlateElems}{192\xspace}
\newcommand{\cbkScaleSmall}{20\xspace}
\newcommand{\cbkScaleLarge}{134\xspace}
\newcommand{\cbkScaleNeSmall}{12\xspace}
\newcommand{\cbkScaleNeLarge}{1728\xspace}
\newcommand{\cbkLoopSlope}{0.98\xspace}
\newcommand{\cbkBatchSlope}{0.83\xspace}
\newcommand{\cbkNewtonN}{1\,230\xspace}
\newcommand{\cbkNewtonNE}{560\xspace}
\newcommand{\cbkNewtonDefault}{13.4\xspace}
\newcommand{\cbkNewtonBatched}{0.3\xspace}
\newcommand{\cbkNewtonRatio}{47\xspace}
\newcommand{\cbkNewtonDiff}{$7.0\times10^{-16}$\xspace}
\newcommand{\cbkNewtonMax}{0.63\xspace}
\newcommand{\cbkAgreeWorst}{$9.6\times10^{-15}$\xspace}
\newcommand{\cbkKFQuad}{138\xspace}
\newcommand{\cbkKTQuad}{88\xspace}
\newcommand{\cbkSFQuad}{135\xspace}
\newcommand{\cbkNEQuad}{1\,600\xspace}
\newcommand{\cbkNDQuad}{3\,362\xspace}
\newcommand{\cbkKFHexEight}{51\xspace}
\newcommand{\cbkKTHexEight}{30\xspace}
\newcommand{\cbkSFHexEight}{54\xspace}
\newcommand{\cbkNEHexEight}{1\,536\xspace}
\newcommand{\cbkNDHexEight}{6\,075\xspace}
\newcommand{\cbkKFHexTwenty}{24\xspace}
\newcommand{\cbkKTHexTwenty}{16\xspace}
\newcommand{\cbkSFHexTwenty}{22\xspace}
\newcommand{\cbkNEHexTwenty}{432\xspace}
\newcommand{\cbkNDHexTwenty}{6\,951\xspace}
\newcommand{\cbkKFTet}{24\xspace}
\newcommand{\cbkKTTet}{55\xspace}
\newcommand{\cbkSFTet}{29\xspace}
\newcommand{\cbkNETet}{768\xspace}
\newcommand{\cbkNDTet}{4\,131\xspace}
\newcommand{\cbkMachine}{Linux 6.8.0-138-generic, 2 CPUs, NumPy 2.2.6\xspace}

% generated by scripts/ch07_assets.py -- do not edit
\newcommand{\ctcPenN}{841\xspace}
\newcommand{\ctcPenExact}{$0$ (to rounding)\xspace}
\newcommand{\ctcPenAtOne}{0.024\xspace}
\newcommand{\ctcPenAtSix}{$2.8\times10^{-8}$\xspace}
\newcommand{\ctcPenAtEight}{$2.8\times10^{-10}$\xspace}
\newcommand{\ctcCondAtEight}{$2.6\times10^{9}$\xspace}
\newcommand{\ctcCondExact}{38\xspace}
\newcommand{\ctcRelForceSix}{$2.2\times10^{-7}$\xspace}
\newcommand{\ctcRelForceOne}{0.15\xspace}
\newcommand{\ctcActiveThirty}{172\xspace}
\newcommand{\ctcObsN}{1521\xspace}
\newcommand{\ctcObsActive}{295\xspace}
\newcommand{\ctcObsIt}{12\xspace}
\newcommand{\ctcObsItSaddle}{12\xspace}
\newcommand{\ctcObsGapMin}{$0$\xspace}
\newcommand{\ctcObsLamMin}{$0$\xspace}
\newcommand{\ctcObsComp}{$0$\xspace}
\newcommand{\ctcObsRS}{$6.7\times10^{-16}$\xspace}
\newcommand{\ctcTraceFirst}{892\xspace}
\newcommand{\ctcTraceFinal}{295\xspace}
\newcommand{\ctcOffMax}{0.0\xspace}
\newcommand{\ctcMinEig}{$3.7\times10^{-2}$\xspace}
\newcommand{\ctcAgreeN}{484\xspace}
\newcommand{\ctcAgreeRedMax}{$8.9\times10^{-16}$\xspace}
\newcommand{\ctcAgreeSadMax}{$8.1\times10^{-16}$\xspace}
\newcommand{\ctcAgreeLemkeMax}{$6.7\times10^{-16}$\xspace}
\newcommand{\ctcAgreePgsMax}{$3.7\times10^{-10}$\xspace}
\newcommand{\ctcAgreeActMin}{20\xspace}
\newcommand{\ctcAgreeActMax}{140\xspace}
\newcommand{\ctcOrdN}{196\xspace}
\newcommand{\ctcOrdTwoLast}{0.92\xspace}
\newcommand{\ctcOrdBdfLast}{2.03\xspace}
\newcommand{\ctcOrdBdfFirst}{1.96\xspace}
\newcommand{\ctcOrdRatioMid}{29\xspace}
\newcommand{\ctcOrdRatioFine}{63\xspace}
\newcommand{\ctcOrdRatioCoarse}{14\xspace}
\newcommand{\ctcImpBdfMaxVcoarse}{1.0000\xspace}
\newcommand{\ctcImpBdfEndVcoarse}{0.9782\xspace}
\newcommand{\ctcImpTwoMaxVcoarse}{1.0000\xspace}
\newcommand{\ctcImpTwoEndVcoarse}{0.8766\xspace}
\newcommand{\ctcImpBdfForceVcoarse}{2.0\xspace}
\newcommand{\ctcImpTwoForceVcoarse}{1.2\xspace}
\newcommand{\ctcFreeBdfVcoarse}{0.4672\xspace}
\newcommand{\ctcFreeTwoVcoarse}{0.0768\xspace}
\newcommand{\ctcImpBdfMaxCoarse}{1.0030\xspace}
\newcommand{\ctcImpBdfEndCoarse}{0.9748\xspace}
\newcommand{\ctcImpTwoMaxCoarse}{1.0000\xspace}
\newcommand{\ctcImpTwoEndCoarse}{0.9126\xspace}
\newcommand{\ctcImpBdfForceCoarse}{7.8\xspace}
\newcommand{\ctcImpTwoForceCoarse}{3.4\xspace}
\newcommand{\ctcFreeBdfCoarse}{0.7216\xspace}
\newcommand{\ctcFreeTwoCoarse}{0.1497\xspace}
\newcommand{\ctcImpBdfMaxFine}{1.0060\xspace}
\newcommand{\ctcImpBdfEndFine}{0.9759\xspace}
\newcommand{\ctcImpTwoMaxFine}{1.0000\xspace}
\newcommand{\ctcImpTwoEndFine}{0.9378\xspace}
\newcommand{\ctcImpBdfForceFine}{19.7\xspace}
\newcommand{\ctcImpTwoForceFine}{6.5\xspace}
\newcommand{\ctcFreeBdfFine}{0.9352\xspace}
\newcommand{\ctcFreeTwoFine}{0.2358\xspace}
\newcommand{\ctcImpBdfMaxVfine}{1.0071\xspace}
\newcommand{\ctcImpBdfEndVfine}{0.9748\xspace}
\newcommand{\ctcImpTwoMaxVfine}{1.0000\xspace}
\newcommand{\ctcImpTwoEndVfine}{0.9554\xspace}
\newcommand{\ctcImpBdfForceVfine}{36.9\xspace}
\newcommand{\ctcImpTwoForceVfine}{12.9\xspace}
\newcommand{\ctcFreeBdfVfine}{0.9893\xspace}
\newcommand{\ctcFreeTwoVfine}{0.3208\xspace}
\newcommand{\ctcScaleDenseN}{3481\xspace}
\newcommand{\ctcScaleDenseT}{3.69\xspace}
\newcommand{\ctcScaleDenseRed}{0.14\xspace}
\newcommand{\ctcScaleDenseSad}{0.19\xspace}
\newcommand{\ctcScaleMaxN}{9801\xspace}
\newcommand{\ctcScaleMaxRed}{0.65\xspace}
\newcommand{\ctcScaleMaxSad}{0.98\xspace}
\newcommand{\ctcScaleMaxIt}{23\xspace}
\newcommand{\ctcScaleDiff}{$1.3\times10^{-15}$\xspace}
\newcommand{\ctcMachine}{Linux 6.8.0-138-generic, 2 CPUs, NumPy 2.2.6\xspace}

% generated by scripts/ch08_assets.py -- do not edit
\newcommand{\cpdNSnap}{60\xspace}
\newcommand{\cpdNTest}{40\xspace}
\newcommand{\cpdNFree}{240\xspace}
\newcommand{\cpdSvDecade}{211\xspace}
\newcommand{\cpdModesForFivePct}{6\xspace}
\newcommand{\cpdModesForOnePct}{12\xspace}
\newcommand{\cpdGalerkinErrTwelve}{$8.0\times10^{-3}$\xspace}
\newcommand{\cpdGalerkinErrLast}{$2.1\times10^{-3}$\xspace}
\newcommand{\cpdProjErrLast}{$1.8\times10^{-3}$\xspace}
\newcommand{\cpdSvLast}{$1.8\times10^{-14}$\xspace}
\newcommand{\cpdSvTwelve}{$1.0\times10^{-4}$\xspace}
\newcommand{\cpdGalerkinOverProj}{1.34\xspace}
\newcommand{\cpdMassWins}{21\xspace}
\newcommand{\cpdMassCompared}{24\xspace}
\newcommand{\cpdMassErrTwelve}{$7.7\times10^{-3}$\xspace}
\newcommand{\cpdGalerkinErrFour}{$5.8\times10^{-2}$\xspace}
\newcommand{\cpdGalerkinErrSix}{$3.5\times10^{-2}$\xspace}
\newcommand{\cpdProjErrSix}{$3.4\times10^{-2}$\xspace}
\newcommand{\cpdGalerkinErrEight}{$1.9\times10^{-2}$\xspace}
\newcommand{\cpdMassErrSix}{$3.6\times10^{-2}$\xspace}
\newcommand{\cpdMassOrth}{$1.8\times10^{-15}$\xspace}
\newcommand{\cpdEucOrth}{$2.9\times10^{-15}$\xspace}
\newcommand{\cpdSnapOlhsB}{$4.3\times10^{-3}$\xspace}
\newcommand{\cpdSnapRandB}{$1.4\times10^{-2}$\xspace}
\newcommand{\cpdSnapOlhsA}{$1.6\times10^{-2}$\xspace}
\newcommand{\cpdSnapRandA}{$4.5\times10^{-2}$\xspace}
\newcommand{\cpdOlhsBetterCount}{6\xspace}
\newcommand{\cpdSnapCases}{7\xspace}
\newcommand{\cpdModalFOne}{0.09\xspace}
\newcommand{\cpdModalErrOneAtSix}{$2.7\times10^{-4}$\xspace}
\newcommand{\cpdModalErrFiveAtTwelve}{$3.9\times10^{-2}$\xspace}
\newcommand{\cpdModalErrThreeAtSix}{$1.8\times10^{-2}$\xspace}
\newcommand{\cpdCostSpeedAffine}{38.5\xspace}
\newcommand{\cpdCostSpeedReasm}{252\xspace}
\newcommand{\cpdCostSpeedAffineBig}{2533\xspace}
\newcommand{\cpdCostSpeedReasmBig}{4812\xspace}
\newcommand{\cpdBreakEven}{3\xspace}
\newcommand{\cpdBreakEvenBig}{8\xspace}
\newcommand{\cpdBreakEvenAffBig}{16\xspace}
\newcommand{\cpdBigNf}{1200\xspace}
\newcommand{\cpdOnlineUs}{17\xspace}
\newcommand{\cpdMachine}{Linux 6.8.0-138-generic, 2 CPUs, NumPy 2.2.6\xspace}

% generated by scripts/ch09_assets.py -- do not edit
\newcommand{\cbtNDof}{48\xspace}
\newcommand{\cbtNState}{96\xspace}
\newcommand{\cbtBandMax}{8\xspace}
\newcommand{\cbtFOne}{83.8\xspace}
\newcommand{\cbtNPre}{15\xspace}
\newcommand{\cbtSoarScaledCols}{4\xspace}
\newcommand{\cbtSoarRawCols}{1\xspace}
\newcommand{\cbtStaticNorm}{$1.1\times10^{-6}$\xspace}
\newcommand{\cbtPreErrMax}{$1.9\times10^{-7}$\xspace}
\newcommand{\cbtPreErrMean}{$2.7\times10^{-4}$\xspace}
\newcommand{\cbtHsvRatioTen}{$5.9\times10^{-6}$\xspace}
\newcommand{\cbtHsvRatioTwenty}{$2.0\times10^{-10}$\xspace}
\newcommand{\cbtBoundHolds}{13\xspace}
\newcommand{\cbtBoundCases}{14\xspace}
\newcommand{\cbtBtStable}{14\xspace}
\newcommand{\cbtMissOrder}{26\xspace}
\newcommand{\cbtMissErr}{$2.0\times10^{-19}$\xspace}
\newcommand{\cbtMissBound}{$5.4\times10^{-20}$\xspace}
\newcommand{\cbtMissRel}{$1.9\times10^{-14}$\xspace}
\newcommand{\cbtBoundRatioMin}{0.04\xspace}
\newcommand{\cbtBoundRatioMax}{0.84\xspace}
\newcommand{\cbtBtErrEight}{$6.4\times10^{-6}$\xspace}
\newcommand{\cbtBtBoundEight}{$3.8\times10^{-10}$\xspace}
\newcommand{\cbtBtErrIntEight}{$7.3\times10^{-11}$\xspace}
\newcommand{\cbtRawErrMax}{$3.4\times10^{-2}$\xspace}
\newcommand{\cbtRawStable}{yes\xspace}
\newcommand{\cbtRawHsvOne}{$6.2\times10^{-3}$\xspace}
\newcommand{\cbtPreHsvOne}{$1.0\times10^{-5}$\xspace}
\newcommand{\cbtModalMaxEight}{$1.0\times10^{-5}$\xspace}
\newcommand{\cbtModalMeanEight}{$1.5\times10^{-2}$\xspace}
\newcommand{\cbtModalMaxTwelve}{$3.0\times10^{-6}$\xspace}
\newcommand{\cbtModalMeanTwelve}{$4.5\times10^{-3}$\xspace}
\newcommand{\cbtBtMaxEight}{$6.4\times10^{-6}$\xspace}
\newcommand{\cbtBtMeanEight}{$9.4\times10^{-3}$\xspace}
\newcommand{\cbtBtMaxTwelve}{$4.1\times10^{-7}$\xspace}
\newcommand{\cbtBtMeanTwelve}{$1.8\times10^{-4}$\xspace}
\newcommand{\cbtKryOneMaxEight}{$6.7\times10^{-8}$\xspace}
\newcommand{\cbtKryOneMeanEight}{$9.0\times10^{-7}$\xspace}
\newcommand{\cbtKryOneMaxTwelve}{$2.5\times10^{-8}$\xspace}
\newcommand{\cbtKryOneMeanTwelve}{$1.9\times10^{-9}$\xspace}
\newcommand{\cbtKryTwoMaxEight}{$1.3\times10^{-2}$\xspace}
\newcommand{\cbtKryTwoMeanEight}{$5.6\times10^{-1}$\xspace}
\newcommand{\cbtKryTwoMaxTwelve}{$1.9\times10^{-4}$\xspace}
\newcommand{\cbtKryTwoMeanTwelve}{$1.9\times10^{-3}$\xspace}
\newcommand{\cbtSoarMaxEight}{$1.5\times10^{-4}$\xspace}
\newcommand{\cbtSoarMeanEight}{$1.5\times10^{-3}$\xspace}
\newcommand{\cbtSoarMaxTwelve}{$3.3\times10^{-8}$\xspace}
\newcommand{\cbtSoarMeanTwelve}{$4.6\times10^{-7}$\xspace}
\newcommand{\cbtFarBt}{$6.1\times10^{-2}$\xspace}
\newcommand{\cbtNearBt}{$9.4\times10^{-3}$\xspace}
\newcommand{\cbtFarKryOne}{$5.3\times10^{-2}$\xspace}
\newcommand{\cbtNearKryOne}{$8.9\times10^{-7}$\xspace}
\newcommand{\cbtFarKryTwo}{$1.4\times10^{0}$\xspace}
\newcommand{\cbtNearKryTwo}{$5.6\times10^{-1}$\xspace}
\newcommand{\cbtFarSoar}{$1.2\times10^{0}$\xspace}
\newcommand{\cbtNearSoar}{$1.4\times10^{-3}$\xspace}
\newcommand{\cbtPassFull}{$4.5\times10^{-4}$\xspace}
\newcommand{\cbtNegCount}{0\xspace}
\newcommand{\cbtRowCount}{15\xspace}
\newcommand{\cbtUnstCount}{0\xspace}
\newcommand{\cbtMachine}{Linux 6.8.0-138-generic, 2 CPUs, NumPy 2.2.6\xspace}

% generated by scripts/ch10_assets.py -- do not edit
\newcommand{\cfrNFree}{40\xspace}
\newcommand{\cfrFOne}{83.8\xspace}
\newcommand{\cfrSelFive}{1.00, 7.18\xspace}
\newcommand{\cfrGreedyNearFive}{$4.2\times10^{-4}$\xspace}
\newcommand{\cfrGreedyMedFive}{$3.0\times10^{-6}$\xspace}
\newcommand{\cfrGreedyBandFive}{$4.2\times10^{-4}$\xspace}
\newcommand{\cfrGreedyNearEight}{$3.2\times10^{-9}$\xspace}
\newcommand{\cfrGreedyMedEight}{$8.5\times10^{-12}$\xspace}
\newcommand{\cfrGreedyBandEight}{$7.4\times10^{-11}$\xspace}
\newcommand{\cfrUniformNearFive}{$1.3\times10^{-4}$\xspace}
\newcommand{\cfrUniformMedFive}{$1.6\times10^{-5}$\xspace}
\newcommand{\cfrUniformBandFive}{$1.3\times10^{-4}$\xspace}
\newcommand{\cfrUniformNearEight}{$5.1\times10^{-9}$\xspace}
\newcommand{\cfrUniformMedEight}{$2.1\times10^{-11}$\xspace}
\newcommand{\cfrUniformBandEight}{$2.2\times10^{-10}$\xspace}
\newcommand{\cfrModalNearFive}{$3.7\times10^{-2}$\xspace}
\newcommand{\cfrModalMedFive}{$1.9\times10^{-3}$\xspace}
\newcommand{\cfrModalBandFive}{$7.3\times10^{-2}$\xspace}
\newcommand{\cfrModalNearEight}{$1.2\times10^{-2}$\xspace}
\newcommand{\cfrModalMedEight}{$6.4\times10^{-4}$\xspace}
\newcommand{\cfrModalBandEight}{$2.4\times10^{-2}$\xspace}
\newcommand{\cfrNTestNear}{60\xspace}
\newcommand{\cfrNTestBand}{100\xspace}
\newcommand{\cfrGreedyWinsNear}{3\xspace}
\newcommand{\cfrRankCount}{7\xspace}
\newcommand{\cfrGreedyWinsModal}{6\xspace}
\newcommand{\cfrGreedyWinsMed}{4\xspace}
\newcommand{\cfrThirdRes}{17.5\xspace}
\newcommand{\cfrSpearResid}{1.00\xspace}
\newcommand{\cfrSpearHier}{0.86\xspace}
\newcommand{\cfrTopOverlapResid}{30\xspace}
\newcommand{\cfrTopOverlapHier}{22\xspace}
\newcommand{\cfrTopK}{30\xspace}
\newcommand{\cfrNCand}{300\xspace}
\newcommand{\cfrHierUnderMax}{22.13\xspace}
\newcommand{\cfrHierOverMax}{4101.1\xspace}
\newcommand{\cfrScmRefs}{15\xspace}
\newcommand{\cfrLpRefs}{15\xspace}
\newcommand{\cfrScmGap}{1.00\xspace}
\newcommand{\cfrLpGap}{1.00\xspace}
\newcommand{\cfrScmValid}{1\xspace}
\newcommand{\cfrLpValid}{1\xspace}
\newcommand{\cfrScmPositive}{1\xspace}
\newcommand{\cfrLpPositive}{8\xspace}
\newcommand{\cfrNGrid}{200\xspace}
\newcommand{\cfrScmBoundHolds}{1\xspace}
\newcommand{\cfrScmBoundFinite}{1\xspace}
\newcommand{\cfrScmRatioMed}{3654\xspace}
\newcommand{\cfrScmRatioMin}{3654.2\xspace}
\newcommand{\cfrScmRatioMax}{$3.7\times10^{3}$\xspace}
\newcommand{\cfrScmInfNear}{25\xspace}
\newcommand{\cfrScmInfFar}{174\xspace}
\newcommand{\cfrScmNearCount}{25\xspace}
\newcommand{\cfrLpBoundHolds}{8\xspace}
\newcommand{\cfrLpBoundFinite}{8\xspace}
\newcommand{\cfrLpRatioMed}{4606\xspace}
\newcommand{\cfrLpRatioMin}{3655.0\xspace}
\newcommand{\cfrLpRatioMax}{$2.0\times10^{4}$\xspace}
\newcommand{\cfrLpInfNear}{25\xspace}
\newcommand{\cfrLpInfFar}{167\xspace}
\newcommand{\cfrLpNearCount}{25\xspace}
\newcommand{\cfrQueryScmUs}{122\xspace}
\newcommand{\cfrQueryLpUs}{1310\xspace}
\newcommand{\cfrQuerySvdUs}{158\xspace}
\newcommand{\cfrSweepLipMax}{200\xspace}
\newcommand{\cfrSweepLipFinite}{1\xspace}
\newcommand{\cfrSweepLipRatio}{3654.2\xspace}
\newcommand{\cfrSweepLpMax}{80\xspace}
\newcommand{\cfrSweepLpFinite}{50\xspace}
\newcommand{\cfrSweepLpRatio}{1178.4\xspace}
\newcommand{\cfrRvGridA}{$1.8\times10^{-2}$\xspace}
\newcommand{\cfrRvGridB}{$9.3\times10^{-3}$\xspace}
\newcommand{\cfrRvNPts}{6000\xspace}
\newcommand{\cfrRvSig}{2.2295e-02\xspace}
\newcommand{\cfrMachine}{Linux 6.8.0-138-generic, 2 CPUs, NumPy 2.2.6\xspace}

% generated by scripts/ch11_assets.py -- do not edit
\newcommand{\ccmBeamFull}{78\xspace}
\newcommand{\ccmBeamFOne}{533.25\xspace}
\newcommand{\ccmCbOneSize}{4\xspace}
\newcommand{\ccmCbOneErr}{$4.8\times10^{-1}$\xspace}
\newcommand{\ccmCbFiveSize}{12\xspace}
\newcommand{\ccmCbFiveErr}{$9.3\times10^{-5}$\xspace}
\newcommand{\ccmCbTwentySize}{42\xspace}
\newcommand{\ccmCbTwentyErr}{$1.0\times10^{-7}$\xspace}
\newcommand{\ccmCbAllSize}{78\xspace}
\newcommand{\ccmCbAllErr}{$5.2\times10^{-11}$\xspace}
\newcommand{\ccmCbMonotone}{yes\xspace}
\newcommand{\ccmGuyTwelveErr}{$1.1\times10^{-3}$\xspace}
\newcommand{\ccmGuyTwentyErr}{$2.0\times10^{-4}$\xspace}
\newcommand{\ccmGuyEightErr}{$4.2\times10^{-3}$\xspace}
\newcommand{\ccmGuyOneErrFirst}{$1.1\times10^{-5}$\xspace}
\newcommand{\ccmGuyTwelveHigh}{$1.1\times10^{-3}$\xspace}
\newcommand{\ccmGuyStatic}{$1.4\times10^{-12}$\xspace}
\newcommand{\ccmMemN}{336\xspace}
\newcommand{\ccmMemIface}{21\xspace}
\newcommand{\ccmMemNmTenErr}{$1.7\times10^{-1}$\xspace}
\newcommand{\ccmMemNmFortyErr}{$2.7\times10^{-3}$\xspace}
\newcommand{\ccmMemNmTenSize}{31\xspace}
\newcommand{\ccmMemNmFortySize}{61\xspace}
\newcommand{\ccmPsiCols}{21\xspace}
\newcommand{\ccmCcTwoErr}{$1.5\times10^{-1}$\xspace}
\newcommand{\ccmGsTwoErr}{$1.5\times10^{-1}$\xspace}
\newcommand{\ccmCcFourErr}{$4.0\times10^{-2}$\xspace}
\newcommand{\ccmGsFourErr}{$4.0\times10^{-2}$\xspace}
\newcommand{\ccmCcEightErr}{$2.7\times10^{-3}$\xspace}
\newcommand{\ccmGsEightErr}{$2.7\times10^{-3}$\xspace}
\newcommand{\ccmCcTwelveErr}{$2.7\times10^{-3}$\xspace}
\newcommand{\ccmGsTwelveErr}{$2.7\times10^{-3}$\xspace}
\newcommand{\ccmCcAllErr}{$2.7\times10^{-3}$\xspace}
\newcommand{\ccmGsAllErr}{$2.7\times10^{-3}$\xspace}
\newcommand{\ccmCcFourMac}{0.0000\xspace}
\newcommand{\ccmGsFourMac}{0.0000\xspace}
\newcommand{\ccmCcEightMac}{0.9978\xspace}
\newcommand{\ccmGsEightMac}{0.9978\xspace}
\newcommand{\ccmCcSizeFour}{44\xspace}
\newcommand{\ccmCcSavedDofs}{17\xspace}
\newcommand{\ccmPlainErr}{$2.7\times10^{-3}$\xspace}
\newcommand{\ccmFreeModalStatic}{$8.5\times10^{-1}$\xspace}
\newcommand{\ccmFreeRubinStatic}{$1.4\times10^{-16}$\xspace}
\newcommand{\ccmFreeMacStatic}{$8.5\times10^{-1}$\xspace}
\newcommand{\ccmFreeCbStatic}{$1.4\times10^{-16}$\xspace}
\newcommand{\ccmFreeModalStaticTen}{$5.3\times10^{-1}$\xspace}
\newcommand{\ccmFreeMacStaticTen}{$5.3\times10^{-1}$\xspace}
\newcommand{\ccmFreeRubinTopThree}{$8.4\times10^{-1}$\xspace}
\newcommand{\ccmFreeCbTopThree}{$8.5\times10^{-1}$\xspace}
\newcommand{\ccmFreeRubinTopTen}{$4.4\times10^{-1}$\xspace}
\newcommand{\ccmFreeCbTopTen}{$3.4\times10^{-1}$\xspace}
\newcommand{\ccmFreeRubinLow}{$3.2\times10^{-14}$\xspace}
\newcommand{\ccmFreeMacLow}{$3.2\times10^{-14}$\xspace}
\newcommand{\ccmFreeCbLowThree}{$1.4\times10^{-2}$\xspace}
\newcommand{\ccmCompFull}{144\xspace}
\newcommand{\ccmCompErrThree}{$5.3\times10^{-2}$\xspace}
\newcommand{\ccmCompErrSix}{$2.1\times10^{-3}$\xspace}
\newcommand{\ccmCompErrTwelve}{$2.8\times10^{-4}$\xspace}
\newcommand{\ccmCompRedTwelve}{24\xspace}
\newcommand{\ccmCompTie}{$2.2\times10^{-16}$\xspace}
\newcommand{\ccmCompConv}{all\xspace}
\newcommand{\ccmCompReuseErr}{$4.0\times10^{-4}$\xspace}
\newcommand{\ccmCompReuseConv}{yes\xspace}
\newcommand{\ccmCompMono}{yes\xspace}
\newcommand{\ccmMachine}{Linux 6.8.0-138-generic, 2 CPUs, NumPy 2.2.6\xspace}

% Measured by the author on his own machine (12 CPUs, Windows) with fea_engine/tests/sparse_contact_timing.py, 2026-10-16.
% Hand-entered from the printed output of that script; to be re-run and regenerated at the pinned release tag.
\newcommand{\ctcAuthN}{39{,}601\xspace}
\newcommand{\ctcAuthActive}{9{,}425\xspace}
\newcommand{\ctcAuthIt}{30\xspace}
\newcommand{\ctcAuthRed}{4.5\xspace}
\newcommand{\ctcAuthSad}{7.6\xspace}
\newcommand{\ctcAuthDiff}{$7.4\times10^{-15}$\xspace}
\newcommand{\ctcAuthDenseGB}{25\xspace}



\title{Architecting Computational Mechanics\\[1ex]{\Large From-Scratch Finite Element Analysis\\and Reduced-Order Modeling in Python}}
\author{Abhijeet Kumar}
\date{Draft 1 --- companion to \texttt{fea\_engine} 1.0.1 and \texttt{rom\_engine} 0.2.0}

\begin{document}
\frontmatter
\maketitle
\chapter*{Reading guide}
This book explains how the two packages of the companion repository, \fea{} and \rom{}, are built, and verifies them. It is a companion to a pinned release of the repository: every listing, number and code anchor refers to that release. Code is cited as \anchor{importable.name} and theory as numbered references. Chapters that are not yet written are marked as such; Chapter~\ref{ch:contact} is written in full and is the template for the others.
\tableofcontents
\listoffigures
\listoftables

\mainmatter
\part{Architectural foundations}
\chapter{The design of a decoupled scientific engine}\label{ch:arch}

\begin{tcolorbox}[colback=gray!5,colframe=gray!50,title={In this chapter}]
The two packages of this book, \fea{} and \rom{}, are meant to grow without rewriting what exists: a new element, material, geometry, load, solution strategy or linear-algebra backend should be added by a subclass or a registry entry. This chapter describes the structure that makes this possible and tests the claims it rests on. It shows the flow from material to post-processing, the small interface that the solvers use to talk to elements, a new element added in about twenty lines to a solver that was not changed, an interface test of all 34 registered elements, the rule that \rom{} never imports \fea{}, and the backend dispatcher, including a cost it has on sparse systems that the book measured and the package documentation does not mention.
\end{tcolorbox}

\section{One system, many problems}\label{sec:ar-system}

Whatever the physical problem is (a membrane, a beam, a shell with contact), the solvers of \fea{} work on one algebraic form,
\begin{equation}\label{eq:ar-system}
\bm M\ddot{\bm u}+\bm C\dot{\bm u}+\bm K\bm u=\bm F ,
\end{equation}
or, for nonlinear problems, on the residual $\bm f_{int}(\bm u)-\bm F$ and its tangent (Chapter~\ref{ch:nonlin}). The vectors and matrices are plain NumPy or SciPy arrays. Everything that is specific to a problem is pushed into the construction of $\bm M,\bm C,\bm K$ (or of $\bm f_{int}$ and $\bm K_T$), and everything that is specific to a solution method is in the solvers. The design question of this chapter is where the boundary between the two sits, and how narrow it can be made.

\section{The flow and the layers}\label{sec:ar-flow}

\Cref{fig:ar-flow} shows the flow of a typical analysis. A \emph{material} gives the constitutive matrix (for example \code{D_plane_stress(Material(...))}); a \emph{mesh} gives nodes and element connectivity (\code{Mesh}, or \code{MultiBlockMesh} for several element topologies, Chapter~\ref{ch:mesh}); an \emph{element} formulation turns coordinates and a material into element matrices or forces; the system \anchor{fea_engine.solver.FESystem} assembles them, applies boundary conditions and loads, and calls a \emph{solver}; \emph{post-processing} recovers stress, strain and fields from the result. A layer knows only the layer next to it: the mesh does not know which element will use it, and the solver does not know which element it is solving for.

\begin{figure}[ht]
\centering
\begin{tikzpicture}[node distance=1.0cm and 0.75cm, box/.style={draw,rounded corners,minimum height=0.9cm,minimum width=2.15cm,align=center,font=\small,fill=gray!8}, arr/.style={-{Stealth},thick}]
\node[box] (mat) {material\\\code{Material}, $\bm D$};
\node[box,right=of mat] (mesh) {mesh\\\code{Mesh}};
\node[box,right=of mesh] (el) {element\\\code{Element}};
\node[box,right=of el] (sys) {system\\\code{FESystem}};
\node[box,below=1.0cm of sys] (sol) {solver\\static, modal,\\transient, Newton};
\node[box,left=of sol] (post) {post-processing\\stress, fields, export};
\node[box,left=of post,fill=blue!8] (rom) {\code{rom_engine}\\plain arrays};
\draw[arr] (mat)--(mesh); \draw[arr] (mesh)--(el); \draw[arr] (el)--(sys); \draw[arr] (sys)--(sol); \draw[arr] (sol)--(post);
\draw[arr,dashed] (sys.south west) .. controls +(-0.3,-0.8) and +(0.3,0.8) .. (rom.north east) node[midway,above,sloped,font=\scriptsize]{$\bm M,\bm K,\bm F$ as arrays};
\end{tikzpicture}
\caption{The flow of an analysis in \fea{}. The dashed arrow is the only link to \rom{}: matrices and vectors as plain arrays, never an import.}\label{fig:ar-flow}
\end{figure}

\section{The element interface}\label{sec:ar-interface}

The narrow boundary is the base class \anchor{fea_engine.elements.base.Element}. An element declares four class attributes, \code{n_nodes}, \code{dofs_per_node}, \code{dim} (the dimension of the natural-coordinate space) and \code{gauss_order}, and provides the shape functions in the natural coordinates, \code{shape_and_derivs}. From these the base class implements the Gauss loop \emph{once}: if the element supplies a strain-displacement matrix \code{B_matrix}, the generic \code{stiffness} computes $\int\bm B^{T}\bm D\bm B$ (\Cref{eq:el-K}) and the generic \code{mass} computes $\int\rho\bm N^{T}\bm N$. An element that needs more overrides the method it needs: the Mindlin plate overrides \code{stiffness} for selective reduced integration, and the Euler--Bernoulli beam uses a closed form. Nonlinear elements provide the pair \code{internal_force} and \code{tangent_stiffness} instead (or as well). Optional capabilities are declared by attributes and methods the solvers look for: \code{supports_batch} with the batched kernels of Chapter~\ref{ch:batch}, \code{state} and \code{commit_state} for history variables (Chapter~\ref{ch:nonlin}), and \code{lumped_mass} for the explicit method (Chapter~\ref{ch:dyn}).

The contact elements of \fea{} are in the same registry but have no interpolation: they define force and tangent from a gap function and do not implement \code{shape_and_derivs}, which raises \code{NotImplementedError} for them. The interface is therefore not a single fixed set of methods, but ``what the solver calls'': the assembly asks the element for matrices (linear) or for force and tangent (nonlinear), and the element is free in how it computes them.

\section{Extending the engine: a new element}\label{sec:ar-extend}

The claim of the design is that a new element touches no existing code. \Cref{lst:ar-bar} defines a two-node bar with one displacement per node (a one-dimensional element, $\bm B=\mathrm dN/\mathrm dx$, with the stiffness $\bm D=EA$ as a $1\times1$ matrix) as a subclass, registers it with one line (after \texttt{from fea\_engine.elements import Element, ELEMENT\_REGISTRY, jacobian}), and uses it with the existing assembly, static solver and modal solver.

\begin{lstlisting}[caption={A new element by subclassing, and the one registry line. Nothing in \texttt{fea\_engine} is edited.},label={lst:ar-bar}]
class Bar1D(Element):
    n_nodes, dofs_per_node, dim, gauss_order = 2, 1, 1, 1
    dof_names = ("u",)

    def shape_and_derivs(self, natural_coords):
        (xi,) = natural_coords
        return np.array([(1 - xi) / 2, (1 + xi) / 2]), np.array([[-0.5, 0.5]])

    def B_matrix(self, natural_coords, elem_coords):
        _, dN = self.shape_and_derivs(natural_coords)
        J, detJ = jacobian(dN, elem_coords)
        return np.linalg.solve(J, dN), detJ

ELEMENT_REGISTRY["bar1d"] = Bar1D                      # the one registry line

m = line_mesh(L, n)
s = FESystem(m, ELEMENT_REGISTRY["bar1d"]())
s.assemble_stiffness(np.array([[E * A]]))
s.assemble_mass(np.array([[rho * A]]))
s.fix_dofs([0], [0]); s.add_nodal_force([n], 0, P)
U = s.solve_static(); f, _ = s.solve_modal(n_modes=1)
\end{lstlisting}

The listing was run for 1, 2, 4, 8, 16 and 32 elements ($E=210$ GPa, $A=10^{-4}$ m$^2$, $\rho=7800$ kg/m$^3$, $L=2$ m). The static tip displacement equals $PL/(EA)$ to \cdeExtStatic{} (the linear element is exact for this problem), and the first axial frequency converges to $\sqrt{E/\rho}/(4L)$ with observed order \cdeExtFreqOrder{}, from an error of \cdeExtFreqCoarse{} with one element. The existing solvers ran on the new element unchanged, which is what the design promises. It does not show that every solver works with every new element: a path-dependent or batched element must also meet the conventions of \Cref{sec:ar-interface}.

\section{The six extension points}\label{sec:ar-points}

\Cref{tab:ar-points} lists where each kind of extension enters. In each case the repository's own convention is a dictionary or a base class, so that a new item is added in one place.

\begin{table}[ht]
\centering\small
\begin{tabularx}{\linewidth}{@{}>{\raggedright\arraybackslash}p{0.17\linewidth} >{\raggedright\arraybackslash}p{0.32\linewidth} X@{}}
\toprule
extension & where it enters & what to provide \\
\midrule
material & \anchor{fea_engine.material.CONSTITUTIVE_REGISTRY}; the \code{D_...} functions & a dataclass and a function returning the constitutive matrix (or a state update for inelastic laws) \\
geometry / mesh & \anchor{fea_engine.geometry.shapes.GEOMETRY_REGISTRY} & a generator returning a \code{Mesh} \\
element & \anchor{fea_engine.elements.ELEMENT_REGISTRY} & an \code{Element} subclass (\Cref{lst:ar-bar}) \\
load & \anchor{fea_engine.loads.LoadPattern}, \code{TimeHistoryLoad}, \code{HarmonicLoad}, \code{PSDLoad} & an object that returns a force vector, for a time or a frequency \\
solution strategy & the drivers of \code{nonlinear_solver}, \code{NewtonOptions} & a function that calls assembly and the linear solve \\
backend & \code{backend=} of \code{FESystem}, \code{backend_dispatch} & \code{"scipy"}, \code{"torch"} or \code{"auto"} (\Cref{sec:ar-backend}) \\
\bottomrule
\end{tabularx}
\caption{Where each kind of extension enters \fea{}.}\label{tab:ar-points}
\end{table}

\section{Does the interface hold? A conformance test}\label{sec:ar-conform}

A registry is only useful if every entry honours the interface. The test of this chapter loops over all \cdeRegTotal{} entries of \code{ELEMENT_REGISTRY} and checks, without needing a material, that each is a subclass of \code{Element}, that \code{n_nodes}, \code{dofs_per_node}, \code{dim} and \code{gauss_order} are positive integers, and, for the interpolated elements, that the shape functions at the natural origin sum to one and that the derivative array has the shape $(\dim,n_{nodes})$ (beam and shell elements, whose first nodal values need not be a partition of unity, skip the sum test). For the contact elements, which have no interpolation, it checks that \code{internal_force} and \code{tangent_stiffness} are callable. \Cref{tab:ar-registry} gives the result: all \cdeRegPass{} entries pass. The test checks the interface and not the physics; the physics is the subject of Chapters~\ref{ch:elem}, \ref{ch:nonlin} and \ref{ch:contact}.

\begin{table}[ht]
\centering\small
\begin{tabular}{lrr}
\toprule
family & registered & pass the interface test \\
\midrule
plane stress / strain & 6 & 6 \\
solid & 11 & 11 \\
plate, shell & 4 & 4 \\
scalar & 1 & 1 \\
beam, truss & 7 & 7 \\
contact & 5 & 5 \\
\midrule
total & 34 & 34 \\
\bottomrule
\end{tabular}

\caption{Interface conformance of the registered elements, by family.}\label{tab:ar-registry}
\end{table}

\section{Two packages, one direction of dependence}\label{sec:ar-packages}

\Cref{tab:ar-packages} gives the size of the packages at the release of this book. The rule that \rom{} works on plain arrays and does not depend on \fea{} can be checked in the source: parsing every module of both packages for import statements shows \cdeFeaImportsRom{} imports of \rom{} in \fea{} and \cdeRomImportsFea{} imports of \fea{} in \rom{}. The tests of \rom{} do import \fea{} (in \cdeRomTestsImportFea{} of \cdeRomTestFiles{} test files), because that is where a reduced model is compared with the full model it came from. The consequence is that \rom{} can reduce a model from any finite element code, as long as the code provides $\bm M$, $\bm K$ and the force as arrays, and that the hooks for hyper-reduction (Chapters~\ref{ch:batch} and~\ref{ch:hyper}) take plain callables instead of an \fea{} object.

\begin{table}[ht]
\centering\small
\begin{tabular}{lrrrr}
\toprule
 & source modules & source lines & test files & test lines \\
\midrule
\code{fea_engine} & 62 & 30\,099 & 135 & 27\,703 \\
\code{rom_engine} & 61 & 17\,013 & 85 & 14\,866 \\
\bottomrule
\end{tabular}

\caption{Size of the packages: Python modules and lines (including comments and docstrings) in the source tree and in the tests.}\label{tab:ar-packages}
\end{table}

\section{Backend dispatch and what it costs}\label{sec:ar-backend}

\code{FESystem(..., backend=...)} chooses the linear-algebra backend of \code{solve_static}: \code{"scipy"} (the default), \code{"torch"} (conjugate gradients on CPU or GPU) or \code{"auto"}. The automatic choice is made at solve time by \anchor{fea_engine.backend_dispatch.select_backend}: a free-dof block that is symmetric positive definite (SPD) and larger than \cdeDispThreshold{} dofs is routed to the torch conjugate-gradient path, everything else to SciPy; a non-SPD matrix is never given to the conjugate-gradient solver. The threshold is a documented rule of thumb and not a benchmark result.

\paragraph{Measured cost of the automatic choice.} \Cref{fig:ar-dispatch} compares \code{solve_static} with \code{backend="scipy"} and \code{backend="auto"} for a sparse plane-stress plate with up to \cdeDispMaxDof{} dofs, a range in which \code{auto} always picks the SciPy path. The two give the same answer, but \code{auto} is slower, and by an amount that grows quickly: at the largest size \code{auto} takes \cdeDispExtraFracLarge{}\% more time than \code{scipy}, \cdeDispExtraMs{} ms more in absolute terms, and the extra time grows with the dof count with an observed exponent of \cdeDispSlope{}. The cause is visible in the source. The SPD test of \code{select_backend} converts the sparse matrix to a dense array and runs a dense Cholesky factorisation, which is $O(n^3)$ in time and $O(n^2)$ in memory, and it does so on every solve, even when the answer decides nothing (below the threshold the SciPy path is taken regardless). For a free-dof count of 5\,000, the dense array takes about \cdeDenseFiveK{} GB; for 50\,000 it would take about \cdeDenseFiveK0{} GB (computed as $8n^2$ bytes, not measured). Until the probe works on the sparse matrix, the practical advice is to pass \code{backend="scipy"} for large sparse systems and to use \code{"auto"} only where the system is small.

\begin{figure}[ht]
\centering
\includegraphics[width=0.95\linewidth]{figures/ch01_dispatch.pdf}
\caption{\code{solve_static} time on a sparse plane-stress plate with \code{backend="scipy"} and \code{backend="auto"} (left), and the extra time of \code{auto} (right). Minimum of three runs on \machine{}. The torch path was not run (torch is \cdeHasTorch{} in the sandbox), and none of the sizes exceeds the threshold at which \code{auto} would choose it.}\label{fig:ar-dispatch}
\end{figure}

\begin{table}[ht]
\centering\small
\begin{tabularx}{\linewidth}{@{}>{\raggedright\arraybackslash}p{0.38\linewidth} X@{}}
\toprule
Concept & Code \\
\midrule
system, assembly, boundary conditions, solvers & \anchor{fea_engine.solver.FESystem} \\
element base class and quadrature & \anchor{fea_engine.elements.base.Element}, \anchor{fea_engine.elements.base.gauss_product} \\
registries & \anchor{fea_engine.elements.ELEMENT_REGISTRY}, \anchor{fea_engine.geometry.shapes.GEOMETRY_REGISTRY}, \anchor{fea_engine.material.CONSTITUTIVE_REGISTRY} \\
dispatcher & \anchor{fea_engine.backend_dispatch.select_backend} \\
\bottomrule
\end{tabularx}
\caption{Where the concepts of this chapter live in the repository (release \booktag).}\label{tab:ar-map}
\end{table}

\begin{limitsbox}
\begin{itemize}
\item \textbf{The interface test is shallow.} It checks attributes, shape functions and the presence of the force and tangent methods, not that an element is correct (\Cref{tab:ar-registry}).
\item \textbf{Extension is easy for the common case.} A new element with the generic Gauss loop is about twenty lines. Elements with history variables, mixed formulations or batched kernels have to follow the conventions of Chapters~\ref{ch:nonlin} and~\ref{ch:batch}.
\item \textbf{The \code{auto} backend has a large cost on sparse systems} (\Cref{fig:ar-dispatch}), because of its dense SPD probe, and the torch branch was not exercised here.
\item \textbf{Line counts} include comments and docstrings, so they measure documentation as well as code.
\item \textbf{One model of extension.} Registries are Python dictionaries: an extension is made in code by the user, and there is no plug-in discovery.
\end{itemize}
\end{limitsbox}

\section{Summary}\label{sec:ar-summary}

\fea{} separates material, mesh, element, system, solver and post-processing, and the only thing a solver needs from an element is the matrix or the force and tangent. The consequence was tested: a new element of twenty lines worked with the existing static and modal solvers with an error of \cdeExtStatic{}; all \cdeRegTotal{} registered elements satisfy the interface test; and \rom{} is independent of \fea{} at the source level. The measured cost of \code{backend="auto"} is an example of why the book tests design claims: the dispatcher works, but not for free.

\section*{Exercises}\label{sec:ar-exercises}
\begin{enumerate}
\item Add a two-node element for a linear spring ($\bm K=k[[1,-1],[-1,1]]$) as in \Cref{lst:ar-bar}, using \code{stiffness} directly instead of \code{B_matrix}. Which of the interface checks of \Cref{sec:ar-conform} does it still need to pass?
\item Extend the conformance test to call \code{stiffness} with a suitable material for every non-contact element and check symmetry and the number of zero eigenvalues (rigid-body modes). Which elements need which material argument?
\item Reproduce \Cref{fig:ar-dispatch} with the dense Cholesky of the probe timed on its own. Does it account for the extra time of \code{auto}?
\item Write a one-line change of \code{_is_spd} that tests a sparse matrix without densifying it (for example with a sparse Cholesky or an \code{eigsh} test of the smallest eigenvalue), and compare the time with \Cref{fig:ar-dispatch}.
\item Parse the source of \rom{} for the modules that mention \fea{} in comments or docstrings. Is any of them a dependency in disguise (a hidden convention)?
\end{enumerate}

\chapterrefs

\chapter{Mesh topologies and structural domain partitioning}\label{ch:mesh}
\stubchapter{Chapter 2}

\part{Full-order physics (\texttt{fea\_engine})}
\chapter{Linear continua and structural elements}\label{ch:elem}

\begin{tcolorbox}[colback=gray!5,colframe=gray!50,title={In this chapter}]
Everything in the later chapters is built from element matrices, so this chapter starts at the root: how a weak form becomes a sparse stiffness matrix and a mass matrix, and which element families \fea{} offers for scalar fields, plane and solid continua, beams, plates and shells. It does not re-derive shape functions. It looks instead at the places where an element can be wrong while looking plausible: locking, integration rules and the choice of mass matrix. Each is followed by a test against an independent reference: patch tests on distorted meshes, pure bending, a slender cantilever of a nearly incompressible material, beam frequencies against closed forms, and a plate against the shell that must reproduce it.\par\smallskip
After it you should be able to say how a weak form becomes a stiffness and a mass matrix, recognise locking or a wrong integration rule from a test with an independent reference, and choose between a consistent and a lumped mass matrix, knowing the frequency error of each.
\end{tcolorbox}

\section{From the weak form to an element matrix}\label{sec:el-weak}

The smallest complete example is the scalar element \anchor{fea_engine.elements.scalar.Quad4Laplace}, which has one unknown per node. For the equation $-\nabla\!\cdot(\kappa\nabla u)=q$ on a domain $\Omega$ with $u=0$ on a part of its boundary, the weak form is: find $u$ such that for every admissible $v$
\begin{equation}\label{eq:el-weak}
\int_\Omega\kappa\,\nabla v\cdot\nabla u\,d\Omega=\int_\Omega v\,q\,d\Omega .
\end{equation}
Inserting the interpolation $u\approx\sum_a N_a u_a$ and the same functions for $v$ turns \eqref{eq:el-weak} into $\bm K\bm u=\bm f$ with
\begin{equation}\label{eq:el-K}
\bm K_e=\int_{\Omega_e}\bm B^{T}\kappa\,\bm B\,d\Omega,\qquad \bm B=\nabla\bm N,\qquad
\bm f_e=\int_{\Omega_e}\bm N^{T}q\,d\Omega .
\end{equation}
The integrals are computed by Gauss quadrature on the reference element, with $d\Omega=\det\bm J\,d\xi$ and $\bm B$ mapped by the Jacobian $\bm J$. The elasticity elements have the same structure with a vector field: $\bm K_e=\int\bm B^{T}\bm D\bm B\,d\Omega$, where $\bm B$ is the strain-displacement matrix and $\bm D$ the elasticity matrix \citep{zienkiewicz2013fem,bathe2014fep}. Everything else in this chapter is a choice of $\bm N$, of $\bm D$, or of the quadrature rule in \eqref{eq:el-K}.

\tikzfigure[0.98\linewidth]{ch03_weak2matrix}{From the weak form to the matrices. The choices that decide whether an element is right are the interpolation, the material matrix and the quadrature rule.}{fig:el-weak2matrix}

\begin{algbox}{Element stiffness and mass by Gauss quadrature}
\begin{pseudo}
\Require node coordinates $\bm x_e$, material matrix $\bm D$ (or $\kappa$), density $\rho$, Gauss rule $\{\bm\xi_g,w_g\}$
\State $\bm K_e\gets\bm 0$;\ $\bm M_e\gets\bm 0$
\For{each Gauss point $g$}
  \State $\bm N,\ \partial\bm N/\partial\bm\xi\gets$ shape functions at $\bm\xi_g$
  \State $\bm J\gets(\partial\bm N/\partial\bm\xi)\,\bm x_e$;\ $\det\bm J$;\ $\bm B\gets\bm J^{-1}\partial\bm N/\partial\bm\xi$ mapped into strain--displacement form
  \State $\bm K_e\gets\bm K_e+w_g\det\bm J\;\bm B^T\bm D\bm B$
  \State $\bm M_e\gets\bm M_e+w_g\det\bm J\;\rho\bm N^T\bm N$
\EndFor
\If{lumped mass requested}
  \State $\bm M_e\gets$ HRZ lumping of $\bm M_e$ \Comment{diagonal, same total mass}
\EndIf
\Ensure $\bm K_e$, $\bm M_e$
\end{pseudo}
\end{algbox}

The rule evaluates $n_g$ integrands of $O(d^2n_{str})$ each, and it sets the rank of $\bm K_e$. It goes wrong when it is too low (hourglass modes, rank deficiency) or too high for the physics (locking, \Cref{sec:el-locking}), and when the Jacobian is not positive. The code is \anchor{fea_engine.elements.base.Element} (\code{stiffness}, \code{mass}) and \code{assemble_lumped_mass}; the checks are \Cref{sec:el-verif} (patch tests, beams) and \Cref{fig:el-mass} (mass choices).

\paragraph{The mass matrix: consistent or lumped.}
The same interpolation gives the consistent mass matrix $\bm M_e=\int\rho\bm N^{T}\bm N\,d\Omega$, which is sparse with the pattern of $\bm K$. A lumped mass matrix is diagonal: \fea{} provides element-by-element HRZ lumping (\code{assemble_lumped_mass}), which keeps the total mass of the element. The diagonal form is what the explicit time integration of Chapter~\ref{ch:dyn} needs, since it makes the solve for the accelerations trivial. The choice has a measurable effect on the eigenvalues. For the lowest eigenvalue of the unit-square membrane, $\lambda_1=2\pi^2$ for a clamped edge, \Cref{fig:el-mass} shows that the consistent mass overestimates it and the lumped mass underestimates it, both with the error falling like $h^{\cecMassConsOrder}$ and $h^{\cecMassLumpOrder}$. On the finest mesh the consistent mass is high by \cecMassConsFine{} and the lumped mass is low by \cecMassLumpFine{}. The signs are the usual ones for these two choices, and an average of the two is a classical way to reduce the error; this chapter does not test that.

\begin{figure}[ht]
\centering
\includegraphics[width=0.55\linewidth]{figures/ch03_mass.pdf}
\caption{Lowest eigenvalue of the clamped unit-square membrane (\code{Quad4Laplace}), consistent against lumped mass: the relative error in absolute value. The consistent-mass eigenvalue is above the exact one and the lumped-mass eigenvalue is below it. Produced by \texttt{scripts/ch03\_assets.py}.}\label{fig:el-mass}
\end{figure}

\section{The element families}\label{sec:el-families}

\Cref{tab:el-elements} lists the elements treated in this chapter with their nodes, degrees of freedom per node and the \code{gauss_order} attribute of the class. Plane and solid elements share one idea: a displacement-based isoparametric element with polynomial interpolation. The plane elements use $\bm D$ for plane stress or plane strain; the 3-D elements use the $6\times6$ isotropic $\bm D$. Linear triangles and tetrahedra have constant strain; the 4-node quadrilateral and the 8-node hexahedron have a bilinear or trilinear field and therefore represent bending badly (\Cref{sec:el-locking}).

\begin{table}[ht]
\centering\small
\begin{tabular}{llrrr}
\toprule
family & element & nodes & dofs per node & \code{gauss\_order} \\
\midrule
scalar & \code{Quad4Laplace} & 4 & 1 & 2 \\
plane & \code{Tri3PlaneStress} & 3 & 2 & 1 \\
plane & \code{Tri6PlaneStress} & 6 & 2 & 2 \\
plane & \code{Quad4PlaneStress} & 4 & 2 & 2 \\
plane & \code{Quad8PlaneStress} & 8 & 2 & 3 \\
solid & \code{Tet4Solid3D} & 4 & 3 & 1 \\
solid & \code{Tet10Solid3D} & 10 & 3 & 2 \\
solid & \code{Hex8Solid3D} & 8 & 3 & 2 \\
solid & \code{Hex8SolidBbar} & 8 & 3 & 2 \\
solid & \code{Hex20Solid3D} & 20 & 3 & 3 \\
beam & \code{Beam2DEulerBernoulli} & 2 & 2 & 1 \\
beam & \code{Beam2DReissner} & 2 & 3 & 2 \\
beam & \code{Beam3DEulerBernoulli} & 2 & 6 & 2 \\
plate/shell & \code{Quad4MindlinPlate} & 4 & 3 & 2 \\
plate/shell & \code{Shell4MITC} & 4 & 6 & 2 \\
\bottomrule
\end{tabular}

\caption{Elements of this chapter. \code{gauss_order} is the attribute of the class as it stands in tag \booktag; its meaning is the number of points per direction for quadrilaterals and hexahedra, and a simplex rule order for triangles and tetrahedra.}\label{tab:el-elements}
\end{table}

\paragraph{Structural elements.}
The beam elements use a one-dimensional mesh. \anchor{fea_engine.elements.beams.Beam2DEulerBernoulli} has two dofs per node (transverse displacement and rotation) and cubic Hermite interpolation. \anchor{fea_engine.elements.beams.Beam2DReissner} adds shear deformation, so it has an axial displacement as well, and \anchor{fea_engine.elements.beams3d.Beam3DEulerBernoulli} has six dofs per node. The corotational versions of these beams belong to Chapter~\ref{ch:nonlin}. \anchor{fea_engine.elements.plates.Quad4MindlinPlate} is a 4-node plate with deflection and two rotations per node, and \anchor{fea_engine.elements.shells.Shell4MITC} is a 4-node shell with membrane, bending and transverse-shear action and six dofs per node, with the assumed transverse shear strain of \citet{dvorkin1984}.

\section{Locking and how the elements avoid it}\label{sec:el-locking}

Locking is an element behaving too stiffly for a deformation mode that its interpolation cannot reproduce without also producing a parasitic strain. It is the one error in this chapter that grows when the problem gets \emph{easier} (a thinner beam, a more incompressible material), which is why it needs a test of its own.

\tikzfigure[0.98\linewidth]{ch03_locking}{The mechanism of shear locking in a bilinear element under bending (schematic). The exact bending mode curves the edges and has no shear strain; the bilinear interpolation keeps the edges straight, which produces a parasitic shear strain and an artificially stiff element.}{fig:el-lockmech}

\paragraph{Bending of continuum elements.}
A 4-node quadrilateral in bending cannot represent the curvature of a beam without shear strain, and the parasitic shear stiffens the beam. \Cref{fig:el-bending} loads a plane-stress beam by a pure end moment and plots the ratio of the computed tip deflection to the exact one (the exact solution is quadratic in the displacements, so it lies in the space of the quadratic elements). The linear elements converge from below: at 4 elements along the beam the ratio is \cecBendQuadFour{} for \code{Quad4PlaneStress} and \cecBendTriFour{} for \code{Tri3PlaneStress}, and at 32 elements it is \cecBendQuadThirtytwo{} and \cecBendTriThirtytwo{}. The observed order of the quadrilateral error is \cecBendQuadOrder, close to the expected second order. The quadratic elements \code{Quad8PlaneStress} and \code{Tri6PlaneStress} reproduce the exact solution up to rounding (\cecBendQuadEightErr{} and \cecBendTriSixErr{} for the tested meshes), as the theory says they must.

\begin{figure}[ht]
\centering
\includegraphics[width=0.6\linewidth]{figures/ch03_bending.pdf}
\caption{Pure bending of a plane-stress beam, tip deflection divided by the exact value, against the number of elements along the beam. The quadratic elements are exact on these meshes. The helper is the one used by \texttt{fea\_engine/tests/test\_benchmarks\_convergence.py}.}\label{fig:el-bending}
\end{figure}

\paragraph{Volumetric locking and mean-dilatation B-bar.}
A nearly incompressible material ($\nu\to0.5$) forces $\operatorname{tr}\bm\varepsilon\to0$ at every Gauss point, which a low-order element satisfies only by not deforming. The mean-dilatation (B-bar) element of \citet{hughes1980bbar}, \anchor{fea_engine.elements.solids.Hex8SolidBbar}, replaces the volumetric part of $\bm B$ by its average over the element, which relaxes the constraint to one condition per element. \Cref{tab:el-locking} bends a slender cantilever ($L/h=20$, square section) with a tip load and compares the tip deflection with the Euler--Bernoulli value. At $\nu=0.3$ the plain 8-node hexahedron gives \cecLockHexPlainA{} of the beam value (shear locking, as in two dimensions), and when $\nu\to0.4999$ it collapses to \cecLockHexPlainB{}, because volumetric locking is added. The B-bar element gives \cecLockHexBbarA{} at $\nu=0.3$ and \cecLockHexBbarB{} at $\nu=0.4999$: it removes the \emph{additional} locking caused by incompressibility and leaves the shear locking untouched, which is the claim of the method and nothing more. The 20-node hexahedron gives \cecLockHexTwentyA{} at $\nu=0.3$ and \cecLockHexTwentyB{} at $\nu=0.4999$ on a mesh with the same node spacing, so it locks little in shear but still shows some volumetric locking.

\begin{table}[ht]
\centering\small
\begin{tabular}{lrrrr}
\toprule
element & mesh & dofs & $\nu=0.3$ & $\nu=0.4999$ \\
\midrule
\code{Hex8Solid3D} & $20\times 4\times 4$ & 1575 & 0.711 & 0.177 \\
\code{Hex8SolidBbar} & $20\times 4\times 4$ & 1575 & 0.729 & 0.745 \\
\code{Hex20Solid3D} & $10\times 2\times 2$ & 963 & 0.989 & 0.905 \\
\bottomrule
\end{tabular}

\caption{Tip deflection of a slender cantilever divided by the Euler--Bernoulli value, three hexahedra and two Poisson ratios. The 20-node mesh has the node spacing of the 8-node mesh. A 3-D beam does not reach exactly 1 even for a perfect element (root and shear effects), so only the differences between the rows are significant.}\label{tab:el-locking}
\end{table}

\paragraph{Shear locking in the plate and its cure.}
\anchor{fea_engine.elements.plates.Quad4MindlinPlate} uses full $2\times2$ Gauss integration for the bending energy and a reduced one-point rule for the transverse shear energy (selective reduced integration, the default \code{'sri'}). With full integration for both, the element locks as the plate gets thin: \Cref{tab:el-plate} gives the tip deflection of a cantilever strip ($h/L=0.01$, $W/L=0.2$) divided by the beam-theory value $PL^3/(3D_pW)$ with the plate rigidity $D_p=Eh^3/(12(1-\nu^2))$. This reference is approximate (it ignores the anticlastic effect and shear), so the converged ratio is \cecPlateSriFine{} and not 1; what the table shows is the difference between the columns. With selective reduced integration the ratio is \cecPlateSriCoarse{} on a mesh of 4 elements along the length. With full integration it is \cecPlateFullCoarse{} on the same mesh and still \cecPlateFullFine{} on 64 elements. The shell \code{Shell4MITC}, which uses a different cure for locking (assumed shear strains tied at the element edges), converges to the same value: on the finest mesh the ratio of the shell to the plate is \cecPlateRatioFS. A flat shell must reproduce the plate, so this agreement is a check of both elements; the membrane response of the shell in this test was exactly zero, as it must be for a flat homogeneous single layer.

\begin{table}[ht]
\centering\small
\begin{tabular}{rrrr}
\toprule
elements along $L$ & \code{Quad4MindlinPlate} (SRI) & same, full integration & \code{Shell4MITC} \\
\midrule
4 & 1.059 & 0.005 & 1.051 \\
8 & 1.081 & 0.021 & 1.069 \\
16 & 1.077 & 0.081 & 1.075 \\
32 & 1.078 & 0.262 & 1.077 \\
64 & 1.078 & 0.602 & 1.078 \\
\bottomrule
\end{tabular}

\caption{Tip deflection of a flat cantilever strip divided by the beam-theory value with plate rigidity (an approximate reference), against the number of elements along the length (a quarter as many across).}\label{tab:el-plate}
\end{table}

\section{Verification}\label{sec:el-verif}

\paragraph{Patch tests.}
An element passes the patch test if a mesh of distorted elements reproduces a linear displacement field exactly when the boundary nodes are given the field, and recovers the constant stress. \Cref{tab:el-patch} gives the largest relative error of the displacement and of the stress for the \cecPatchCount{} continuum elements on distorted meshes; the worst value is \cecPatchWorst{}, which is rounding. The patch test is necessary but not sufficient for convergence, and it says nothing about locking, which is why the tests above exist. The B-bar hexahedron is not in this table: it does not support stress recovery in the nodal form the test uses, so its patch test (displacements only) is the one of \texttt{tests/test\_hex8\_bbar.py}, which passes together with the other three tests of that file (4 passed in the run for this chapter).

\begin{table}[ht]
\centering\small
\begin{tabular}{lrr}
\toprule
element & displacement error & stress error \\
\midrule
\code{Quad4PlaneStress} & $5.5\times10^{-16}$ & $5.2\times10^{-15}$ \\
\code{Quad8PlaneStress} & $1.0\times10^{-15}$ & $7.6\times10^{-15}$ \\
\code{Tri3PlaneStress} & $3.7\times10^{-16}$ & $4.9\times10^{-15}$ \\
\code{Tri6PlaneStress} & $1.8\times10^{-15}$ & $1.7\times10^{-14}$ \\
\code{Hex8Solid3D} & $1.8\times10^{-16}$ & $8.9\times10^{-16}$ \\
\code{Hex20Solid3D} & $1.2\times10^{-15}$ & $2.7\times10^{-15}$ \\
\code{Tet4Solid3D} & $3.6\times10^{-16}$ & $8.9\times10^{-16}$ \\
\code{Tet10Solid3D} & $1.2\times10^{-15}$ & $1.9\times10^{-15}$ \\
\bottomrule
\end{tabular}

\caption{Patch tests on distorted meshes. The test helper is \texttt{\_patch} of \texttt{test\_benchmarks\_convergence.py}.}\label{tab:el-patch}
\end{table}

\paragraph{Beams.}
The Euler--Bernoulli element is exact for an end load: one element reproduces the tip deflection $PL^3/(3EI)$ to \cecEBTipErr, because the cubic Hermite interpolation contains the exact solution. For vibration the cubic Hermite element converges with order \cecEBFreqOrder{} in the lowest cantilever frequency (\cecEBFreqHz{} Hz for the test beam), the expected fourth order. The Reissner beam approaches the Timoshenko tip deflection $PL^3/(3EI)+PL/(GA)$ with order \cecTimoOrder{} (\Cref{fig:el-beams}); shear contributes \cecTimoShearShare{}\% of the tip deflection for this slender test beam, so it is a check of the formulation more than a demonstration of shear effects.

\begin{figure}[ht]
\centering
\includegraphics[width=0.95\linewidth]{figures/ch03_beams.pdf}
\caption{Convergence of the first cantilever frequency of the Euler--Bernoulli beam (left) and of the tip deflection of the Reissner beam (right) with the element size, with the observed order.}\label{fig:el-beams}
\end{figure}

\begin{table}[ht]
\centering\small
\begin{tabularx}{\linewidth}{@{}>{\raggedright\arraybackslash}p{0.44\linewidth} X >{\centering\arraybackslash}p{0.11\linewidth}@{}}
\toprule
Concept & Code & Script calls \\
\midrule
scalar element, weak form, mass & \anchor{fea_engine.elements.scalar.Quad4Laplace} & yes \\
plane and solid elements & \anchor{fea_engine.elements.solids.Quad4PlaneStress}, \anchor{fea_engine.elements.solids.Quad8PlaneStress}, \anchor{fea_engine.elements.solids.Tri3PlaneStress}, \anchor{fea_engine.elements.solids.Tri6PlaneStress}, \anchor{fea_engine.elements.solids.Hex8Solid3D}, \anchor{fea_engine.elements.solids.Hex20Solid3D}, \anchor{fea_engine.elements.solids.Tet4Solid3D}, \anchor{fea_engine.elements.solids.Tet10Solid3D} & yes \\
B-bar hexahedron & \anchor{fea_engine.elements.solids.Hex8SolidBbar} & yes \\
beams & \anchor{fea_engine.elements.beams.Beam2DEulerBernoulli}, \anchor{fea_engine.elements.beams.Beam2DReissner}, \anchor{fea_engine.elements.beams3d.Beam3DEulerBernoulli} & yes \\
plate, shell & \anchor{fea_engine.elements.plates.Quad4MindlinPlate}, \anchor{fea_engine.elements.shells.Shell4MITC} & yes \\
convergence studies & \anchor{fea_engine.convergence.run_study}, \anchor{fea_engine.convergence.observed_order} & yes \\
\bottomrule
\end{tabularx}
\caption{Where the concepts of this chapter live in the repository (tag \booktag). The last column says whether the chapter's script \texttt{scripts/chNN\_assets.py} calls, by name, the code listed in the row: yes (every name), part, or no. Code reached only through another routine counts as no, so the column is a lower bound on what this book exercises.}\label{tab:el-map}
\end{table}

\begin{limitsbox}
\begin{itemize}
\item \textbf{Locking is not removed by the patch test.} All elements of \Cref{tab:el-patch} pass it, and the plain 8-node hexahedron still gives \cecLockHexPlainA{} of the beam value. For slender or nearly incompressible problems, choose the element by the tests of \Cref{sec:el-locking}, not by the patch test.
\item \textbf{B-bar addresses volumetric locking only.} It does not cure shear locking (\Cref{tab:el-locking}), and the 20-node element still showed some volumetric locking at $\nu=0.4999$ on the mesh used.
\item \textbf{Stress recovery} is not available for every element (for example \code{Hex8SolidBbar}); use \code{at='centroid'} where the error message says so.
\item \textbf{The plate reference is approximate.} \Cref{tab:el-plate} compares with a strip of plate rigidity, so only differences between its columns are claimed.
\item \textbf{Not covered:} orthotropic and layered materials, thermal and piezoelectric coupling, shells in large rotation (Chapter~\ref{ch:nonlin} treats the corotational and the director shell only in outline), and meshes of high aspect ratio or strong distortion beyond the patch-test meshes.
\end{itemize}
\end{limitsbox}

\section{Summary}\label{sec:el-summary}

Every element is $\int\bm B^{T}\bm D\bm B$ evaluated by Gauss quadrature, plus a mass matrix that is consistent or lumped; the consistent mass overestimates the membrane eigenvalue and the lumped mass underestimates it. All continuum elements pass the patch test on distorted meshes to rounding, but the low-order ones lock in bending, and nearly incompressible material locks the plain hexahedron to \cecLockHexPlainB{} of the beam value. B-bar fixes the incompressible part only, selective reduced integration keeps the Mindlin plate usable when thin, and the MITC shell agrees with the plate in the flat limit. The beams are exact (Euler--Bernoulli, end load) or converge at the expected orders.

\section*{Exercises}\label{sec:el-exercises}
\begin{enumerate}
\item \emph{(Implementation)} Repeat \Cref{fig:el-mass} with the average of the consistent and the lumped matrix, $\bm M=\tfrac12(\bm M_c+\bm M_l)$. Does the error of $\lambda_1$ cancel to leading order, and what order does it have?
\item \emph{(Implementation)} Run the locking test of \Cref{tab:el-locking} on three meshes, $10\times2\times2$, $20\times4\times4$ and $40\times8\times8$, for the plain and the B-bar hexahedron. Does shear locking go away with refinement, and at what rate?
\item \emph{(Implementation)} Use \code{integration='full'} for the plate on a thicker strip ($h/L=0.1$). Is the locking still visible? Compare with the Timoshenko value.
\item \emph{(Interpretation)} Replace the end load of the cantilever by a distributed load. Is one Euler--Bernoulli element still exact? Why not?
\item \emph{(Implementation)} Distort the patch-test mesh more strongly (move the interior nodes by half an element). Does any element in \Cref{tab:el-patch} fail?
\item \emph{(Derivation)} For a fixed--free bar of one two-node element, find the lowest eigenvalue with the consistent mass matrix $\tfrac{\rho AL}{6}\left[\begin{smallmatrix}2&1\\1&2\end{smallmatrix}\right]$ and with the lumped matrix, and show that their errors against the exact $(\pi/2)^2EA/(\rho AL^2)$ are about $+21.6\,\%$ and $-18.9\,\%$.
\end{enumerate}

\chapterrefs

\chapter{Dynamics: modal, transient, harmonic and random response}\label{ch:dyn}
\stubchapter{Chapter 4}

\chapter{Material and geometric nonlinearity}\label{ch:nonlin}
\stubchapter{Chapter 5}

\chapter{Computational acceleration: vectorised assembly and batched kernels}\label{ch:batch}
\stubchapter{Chapter 6}

\chapter{Exact unilateral contact: sparse active sets and BDF2}\label{ch:contact}

\begin{tcolorbox}[colback=gray!5,colframe=gray!50,title={In this chapter}]
A membrane is pressed down onto a bumpy floor. The obvious way to keep it from passing through is to put a stiff spring at every node that has sunk below the floor, and it works, until the spring gets stiff enough for the answer to be good and the matrix gets too ill-conditioned to solve. Contact is the first place where the solver cannot be a single linear solve, because the set of points in contact is unknown and has to be found. This chapter enforces frictionless contact exactly, with Lagrange multipliers, for models with tens of thousands of unknowns. It measures what the penalty costs, derives the primal-dual active-set method as a semismooth Newton iteration, explains the three sparse routes of \anchor{fea_engine.contact_solvers.solve_contact_static_sparse} and how the solver reports a route that does not suit the problem, and then adds a second-order time scheme (BDF2) for impact. Each claim is followed by the measurement that supports it, with the limits stated beside the results.\par\smallskip
After it you should be able to explain why a penalty parameter is the wrong tool for contact, state the primal-dual active-set method and choose among the sparse routes, and run an impact with BDF2 and read its energy behaviour.
\end{tcolorbox}

\section{The problem}\label{sec:ct-problem}

Take a discretised elastic body with $n$ unknowns, a symmetric positive definite stiffness matrix $\bm K\in\mathbb R^{n\times n}$ and a load vector $\bm f$. A set of $m$ linear constraints $\bm B\bm u\ge\bm c$ says that the body may not pass through an obstacle: row $i$ of $\bm B\in\mathbb R^{m\times n}$ maps the displacements to the gap at one potential contact point, and $c_i$ is the offset of the obstacle there. The contact forces are the Lagrange multipliers $\bm\lambda\in\mathbb R^m$, and the discrete frictionless contact problem is
\begin{equation}\label{eq:ct-kkt}
\bm K\bm u=\bm f+\bm B^{T}\bm\lambda,\qquad
\bm g:=\bm B\bm u-\bm c\;\ge\;0,\qquad
\bm\lambda\;\ge\;0,\qquad
\bm\lambda^{T}\bm g=0 .
\end{equation}
The four conditions say, in turn: equilibrium including the contact forces; no penetration; contact forces only push; and a contact force acts only where the gap is closed. The last condition is what makes the problem nonlinear, and it is the same condition throughout the book: the packages use this convention in \fea{} here and in the reduced models of Chapter~\ref{ch:hyper}~\cite{balajewicz2016contact}. For the formulation in a general setting see \citet{wriggers2006contact} and the survey of \citet{wohlmuth2011contact}.

\begin{sidenote}{Lagrange multiplier and KKT}
A Lagrange multiplier $\lambda_i$ is an extra unknown that stands for the force an inequality exerts when it is binding. Here it is the contact force at node $i$. The four conditions of \eqref{eq:ct-kkt} are the Karush--Kuhn--Tucker (KKT) conditions: equilibrium, feasibility, a sign condition on the force, and complementarity, which says that a force can only act where the gap is closed.
\end{sidenote}

\paragraph{The model problem.}
All the experiments of this chapter use one problem that is small enough to inspect and large enough to need a sparse solver: a membrane, a unit square clamped at its edge, loaded by a uniform downward pressure, and pushed onto a rigid obstacle below it. The membrane is discretised with the scalar element \anchor{fea_engine.elements.scalar.Quad4Laplace} (one unknown per node), $\bm B=\bm I$ (every node may touch the obstacle), and the obstacle is a flat floor with two smooth bumps. \Cref{fig:ct-obstacle} shows the solution on a $40\times40$ mesh. The area where the membrane rests on the obstacle, the contact set, is not known in advance; it is found by the algorithm of \Cref{sec:ct-pdas}.

\begin{figure}[ht]
\centering
\includegraphics[width=0.92\linewidth]{figures/ch07_obstacle.pdf}
\caption{The model problem on a $40\times40$ mesh ($\ctcObsN$ unknowns). Left: displacement with the contact set outlined; right: contact force $\bm\lambda$, which is zero outside the contact set. The solver finds $\ctcObsActive$ active constraints in $\ctcObsIt$ iterations.}\label{fig:ct-obstacle}
\end{figure}

For this solution the contact conditions \eqref{eq:ct-kkt} hold to rounding: the smallest gap is \ctcObsGapMin, the smallest contact force is \ctcObsLamMin, and the largest product $\lambda_i g_i$ is \ctcObsComp.

\section{Why not a penalty?}\label{sec:ct-penalty}

The usual way to avoid the inequality is to let the obstacle push back with a spring: for $\bm B=\bm I$, add the force $k_p\max(0,\,c_i-u_i)$ at every node that has penetrated. Then \eqref{eq:ct-kkt} is replaced by $\bm K\bm u=\bm f+k_p\max(0,\bm c-\bm u)$, an ordinary nonlinear equation, and a Newton iteration on the set $\mathcal A=\{i:u_i<c_i\}$ solves $(\bm K+k_p\bm P_{\mathcal A})\bm u=\bm f+k_p\bm P_{\mathcal A}\bm c$, where $\bm P_{\mathcal A}$ selects the penetrating nodes. The penalty method is simple, and \fea{} keeps it as the contact elements \code{GapContactPenalty} and \code{NodeToSegmentContact2D}. Its cost is a parameter that has no physical meaning.

\paragraph{Try it.} \Cref{lst:ct-penalty} builds the membrane of the model problem on a $30\times30$ mesh, solves it by the penalty method at three stiffnesses, and then solves it exactly. The penalty loop is the Newton iteration just described; \code{kp} is measured in units of the mean diagonal of $\bm K$.

\begin{lstlisting}[caption={The membrane obstacle problem: penalty at three stiffnesses against the exact solution (run with the package on the path).},label={lst:ct-penalty}]
import numpy as np, scipy.sparse as sp, scipy.sparse.linalg as spl
from fea_engine import FESystem, Quad4Laplace, solve_contact_static_sparse
from fea_engine import mesh as meshlib

m = meshlib.rectangle_mesh(1.0, 1.0, 30, 30)
fes = FESystem(m, Quad4Laplace(), sparse=True)
fes.assemble_stiffness(1.0); fes.assemble_lumped_mass(1.0)
X = m.nodes
edge = np.flatnonzero((X[:, 0] < 1e-9) | (X[:, 0] > 1 - 1e-9) |
                      (X[:, 1] < 1e-9) | (X[:, 1] > 1 - 1e-9))
fes.fix_dofs(edge, ["u"])
free = np.asarray(fes.free_dofs)
K = sp.csr_matrix(fes.K)[free][:, free]
M = sp.diags(sp.csr_matrix(fes.M_lumped).diagonal()[free]).tocsr()
f = -10.0 * (M @ np.ones(len(free)))                       # uniform pressure
c = np.where(np.hypot(X[free, 0] - .5, X[free, 1] - .5) < .2, -0.05, -10.0)
B = sp.identity(len(free), format="csr")                   # every node may touch

# naive: a penalty spring at every node that has penetrated, solved by Newton
def penalty(kp, iters=60):
    u = np.zeros(len(free))
    for _ in range(iters):
        P = sp.diags((u < c).astype(float))
        A = (K + kp * P).tocsc()
        un = spl.spsolve(A, f + kp * (P @ c))
        if np.allclose(un, u, atol=1e-14): break
        u = un
    return u, A

kbar = K.diagonal().mean()
for k in (1e2, 1e4, 1e6):
    u, A = penalty(k * kbar)
    cond = np.linalg.cond(A.toarray())
    print(f"penalty kp = {k:.0e} x mean diag: penetration {max(0, (c - u).max()):.1e}, "
          f"condition number {cond:.1e}")

r = solve_contact_static_sparse(K, f, B, c)                # route="auto" -> "reduced"
print(f"exact:  penetration {max(0, -r.gap.min()):.1e}, condition number "
      f"{np.linalg.cond(K.toarray()):.1e}, route {r.route}, {r.iterations} iterations, "
      f"{len(r.active)} active")
\end{lstlisting}

\begin{lstlisting}[numbers=none,language={},keywordstyle={},commentstyle={},frame=single,aboveskip=2pt]
penalty kp = 1e+02 x mean diag: penetration 4.3e-04, condition number 3.4e+03
penalty kp = 1e+04 x mean diag: penetration 4.3e-06, condition number 3.4e+05
penalty kp = 1e+06 x mean diag: penetration 4.3e-08, condition number 3.4e+07
exact:  penetration 0.0e+00, condition number 1.8e+02, route reduced, 2 iterations, 111 active
\end{lstlisting}

Each factor of $100$ in the stiffness lowers the penetration by $100$ and raises the condition number by $100$: the penetration shrinks like $1/k_p$ and the conditioning worsens like $k_p$, and neither reaches the exact values. The exact solve has no such parameter, solves a system with a condition number of $180$, and finds the $111$ contact points in two iterations. \Cref{fig:ct-penalty} repeats the measurement on a finer mesh.

\Cref{fig:ct-penalty} measures that cost on the model problem with $\ctcPenN$ unknowns. The penetration shrinks like $1/k_p$, from $\ctcPenAtOne$ at $k_p=\overline{K_{ii}}$ (the mean diagonal of $\bm K$) to \ctcPenAtEight at $k_p=10^8\,\overline{K_{ii}}$, and the penalty force approaches the exact multipliers (largest relative error $\ctcRelForceOne$ at the smallest $k_p$, \ctcRelForceSix at $k_p=10^6\,\overline{K_{ii}}$). At the same time the condition number of the matrix grows in proportion to $k_p$, reaching \ctcCondAtEight at $k_p=10^8\,\overline{K_{ii}}$, against $\ctcCondExact$ for the matrix $\bm K_{FF}$ of the free unknowns that the exact method solves. A penetration small enough for the application therefore has to be bought with an ill-conditioned system, and for an iterative solver or a single-precision GPU run the price is larger still. The exact method has no such parameter: the penetration is \ctcPenExact.

\begin{figure}[ht]
\centering
\includegraphics[width=0.96\linewidth]{figures/ch07_penalty_vs_exact.pdf}
\caption{Penalty against exact contact on the model problem ($\ctcPenN$ unknowns). Left: the largest penetration falls like $1/k_p$ but never reaches zero; the exact solution has no penetration (drawn at the plotting floor). Right: the condition number of the penalty matrix grows in proportion to $k_p$; the exact route solves a system whose condition number does not depend on any parameter.}\label{fig:ct-penalty}
\end{figure}

\section{Mathematical structure}\label{sec:ct-math}

\paragraph{A minimisation problem.}
Because $\bm K$ is symmetric positive definite, \eqref{eq:ct-kkt} are the Karush--Kuhn--Tucker conditions of the convex quadratic program
\begin{equation}\label{eq:ct-qp}
\min_{\bm u}\ \tfrac12\bm u^T\bm K\bm u-\bm f^T\bm u\quad\text{subject to}\quad \bm B\bm u\ge\bm c,
\end{equation}
with Lagrangian $\mathcal L(\bm u,\bm\lambda)=\tfrac12\bm u^T\bm K\bm u-\bm f^T\bm u-\bm\lambda^T(\bm B\bm u-\bm c)$. The constraints are linear and the objective is strictly convex, so the conditions are necessary and sufficient, and a minimiser exists and is unique whenever the feasible set is not empty~\cite{wriggers2006contact}. This is the energy view of contact: the body takes the shape of least potential energy among the shapes that do not penetrate the obstacle.

\paragraph{The condensed problem.}
Eliminating $\bm u=\bm K^{-1}(\bm f+\bm B^T\bm\lambda)$ gives a problem in the multipliers alone,
\begin{equation}\label{eq:ct-lcp}
\bm g=\bm A\bm\lambda+\bm q\ \ge 0,\qquad \bm\lambda\ge0,\qquad \bm\lambda^T\bm g=0,\qquad \bm A=\bm B\bm K^{-1}\bm B^T,\quad \bm q=\bm B\bm K^{-1}\bm f-\bm c,
\end{equation}
a linear complementarity problem (LCP). The matrix $\bm A$ is positive definite when $\bm B$ has full row rank, and then \eqref{eq:ct-lcp} has exactly one solution~\cite{cottle1992lcp}. The dense solvers of \fea{} work on this form: Lemke's pivoting method~\cite{lemke1965bimatrix}, projected Gauss--Seidel, and an active-set method. They share one limitation. The matrix $\bm A$ is $m\times m$ and dense, and forming it needs $\bm K^{-1}\bm B^T$, an $n\times m$ dense array. Counting both arrays in double precision, the memory is about $8(nm+m^2)$ bytes. For $n=m=\ctcAuthN$, the size of the largest case of \Cref{sec:ct-verif-static}, that is about $\ctcAuthDenseGB$~GB, so the condensed form is a reference and a tool for small problems, not a method for a model of this size.

\paragraph{Complementarity as an equation.}
The three last conditions of \eqref{eq:ct-kkt} are equivalent, for any $c_p>0$, to the single componentwise equation
\begin{equation}\label{eq:ct-nonsmooth}
\bm\lambda-\max\bigl(\bm 0,\ \bm\lambda-c_p\,\bm g\bigr)=\bm 0 .
\end{equation}
Indeed, if $g_i>0$ the maximum is $0$ and the equation forces $\lambda_i=0$; if $\lambda_i>0$ the maximum is $\lambda_i-c_pg_i$ and the equation forces $g_i=0$; and if $g_i<0$ the maximum exceeds $\lambda_i$ and the equation reads $c_pg_i=0$, which is impossible. The function is not differentiable where $\lambda_i=c_pg_i$, but it is piecewise linear, and Newton's method can be applied to it with a generalised derivative. The result is the primal-dual active-set method~\cite{hintermuller2002pdas}, and that equivalence is why the method is fast: it is a Newton method, not a pivoting rule.

\section{The primal-dual active-set method and its three routes}\label{sec:ct-pdas}

\paragraph{The algorithm.}
Let $\mathcal A_k$ be the set of constraints assumed active at iteration $k$. The Newton step for \eqref{eq:ct-nonsmooth} treats the constraints in $\mathcal A_k$ as equalities $(\bm B\bm u)_i=c_i$, sets $\lambda_i=0$ outside, solves the resulting linear system, and then reclassifies each constraint by the sign of $\lambda_i-c_pg_i$.

\begin{algbox}{Primal-dual active set (semismooth Newton)}
\begin{pseudo}
\Require $\bm K,\bm f,\bm B,\bm c$, scale $c_p$, initial active set $\mathcal A_0$ (empty, or a warm start)
\Repeat
  \State solve $\bm K\bm u-\bm B_{\mathcal A}^T\bm\lambda_{\mathcal A}=\bm f$,\ $\bm B_{\mathcal A}\bm u=\bm c_{\mathcal A}$;\ $\lambda_i\gets0$ for $i\notin\mathcal A$ \Comment{\texttt{reduced}, \texttt{saddle} or \texttt{condensed} route}
  \State $\bm g\gets\bm B\bm u-\bm c$
  \State $\mathcal A_{k+1}\gets\{i:\lambda_i-c_pg_i>0\}$
  \If{$\mathcal A_{k+1}$ was seen before} \textbf{report cycling and stop} \EndIf
\Until{$\mathcal A_{k+1}=\mathcal A_k$ and the contact conditions hold to the tolerance}
\Ensure $\bm u$, $\bm\lambda\ge\bm0$, active set
\end{pseudo}
\end{algbox}

Each iteration is one sparse factorisation, of the free block for the \texttt{reduced} route and of the saddle-point matrix for \texttt{saddle}, and a warm start from a nearby problem saves iterations. If $\bm K$ is not an M-matrix there is no finite-termination guarantee, and the solver reports cycling or stalling; the dense \texttt{condensed} route needs about $8(nm+m^2)$ bytes. The code is \anchor{fea_engine.contact_solvers.solve_contact_static_sparse}; \Cref{fig:ct-trace} and \Cref{tab:ct-agree} (agreement with the dense solvers) are the checks.

\tikzfigure[0.98\linewidth]{ch07_complementarity}{Left: the complementarity condition (red) against a penalty law (blue, dashed), which lets the gap become negative. Right: the active-set iteration in four steps.}{fig:ct-idea}

The constant $c_p$ sets the scale on which gap and force are compared; the implementation takes the mean diagonal of $\bm K$ (divided by the mean squared row norm of $\bm B$ for a general $\bm B$), so that the two terms have the same units. Starting from an empty set, the first iterate is the unconstrained solution; every constraint it violates enters the active set, and constraints that were wrongly included leave in later iterations. \Cref{fig:ct-trace} shows this on the model problem: the active set starts with $\ctcTraceFirst$ constraints and shrinks monotonically to the final $\ctcTraceFinal$, and the number of constraints that change status per iteration falls to zero. Both routes have the same number of active constraints at every iteration.

\begin{figure}[ht]
\centering
\includegraphics[width=0.96\linewidth]{figures/ch07_active_trace.pdf}
\caption{The active-set iteration on the model problem with $N=\ctcObsN$ unknowns. Left: the number of active constraints per iteration for the \texttt{reduced} and \texttt{saddle} routes (the curves coincide). Right: the number of constraints that changed status; the last point (plotted at 0.5) is zero changes, which is the stopping test.}\label{fig:ct-trace}
\end{figure}

\paragraph{Finite termination, and where it is not guaranteed.}
For a stiffness matrix that is an M-matrix (positive definite with non-positive off-diagonal entries) the method is known to terminate after finitely many iterations~\cite{hintermuller2002pdas}. The membrane stiffness of this chapter satisfies the condition: on the mesh used for the agreement study ($\ctcAgreeN$ unknowns) the largest off-diagonal entry of $\bm K$ is $\ctcOffMax$ and its smallest eigenvalue is \ctcMinEig, so $\bm K$ is positive definite. The stiffness of an elastic solid generally does not, and for it the package gives \emph{no} guarantee. This is the reason for the failure reports of \Cref{sec:ct-impl}: when the iteration cycles or stalls, the solver says so.

\paragraph{Three routes.}
Step 2 is a sparse linear solve, and how it is done defines the routes of \anchor{fea_engine.contact_solvers.solve_contact_static_sparse}.

\begin{description}[leftmargin=1em]
\item[\texttt{reduced}] applies when every row of $\bm B$ selects one node with a positive weight and no node is selected twice ($\bm B=\bm I$ is the typical case). The active unknowns are fixed to their bounds, $u_i=c_i/b_i$, and only the free block is solved: $\bm K_{FF}\bm u_F=\bm f_F-\bm K_{FA}\bm u_A$. The multipliers follow from the residual, $\lambda_i=(\bm K\bm u-\bm f)_i/b_i$ on the active set. One sparse factorisation of the free block per iteration; this is the fast path.
\item[\texttt{saddle}] applies to a general sparse $\bm B$. Each iteration factorises the saddle-point matrix $\begin{bmatrix}\bm K&-\bm B_{\mathcal A}^T\\ \bm B_{\mathcal A}&\bm 0\end{bmatrix}$. The matrix is larger and indefinite, so it costs more, but it needs nothing from $\bm B$ except sparsity.
\item[\texttt{condensed}] forms $\bm A$ of \eqref{eq:ct-lcp} and calls the dense active-set solver. It is the reference and the last resort, and it is guarded by a memory estimate.
\end{description}
\Cref{tab:ct-routes} summarises when each route applies. \texttt{route="auto"} chooses \texttt{reduced} when $\bm B$ is a node selection and \texttt{saddle} otherwise; it never chooses \texttt{condensed} by itself.

\begin{table}[ht]
\centering\small
\begin{tabularx}{\linewidth}{@{}l X X X@{}}
\toprule
Route & applies to & work per iteration & fails with \\
\midrule
\texttt{reduced} & $\bm B$ selects nodes (distinct columns, positive entries) & sparse factorisation of $\bm K_{FF}$ & \texttt{not\_applicable} for a general $\bm B$; \texttt{singular}; \texttt{cycled}; \texttt{max\_iter} \\
\texttt{saddle} & any sparse $\bm B$ & sparse factorisation of the saddle matrix & \texttt{singular} (dependent active rows); \texttt{cycled}; \texttt{max\_iter} \\
\texttt{condensed} & small $n$, $m$ & dense $\bm K^{-1}\bm B^T$ and $\bm A$, $\approx 8(nm+m^2)$ bytes & \texttt{kkt}; \texttt{memory} above \texttt{max\_dense\_bytes} \\
\bottomrule
\end{tabularx}
\caption{The three routes of the sparse contact solver and the failure reasons each can report.}\label{tab:ct-routes}
\end{table}

\section{Implementation}\label{sec:ct-impl}

\paragraph{Design decisions.}
The solver takes plain arrays, $\bm K$, $\bm f$, $\bm B$, $\bm c$, and no mesh or element object, for the reason of Chapter~\ref{ch:arch}: the same function serves the full-order model here and, with other matrices, the reduced models of Chapter~\ref{ch:hyper}. Three rules shape the interface.
\begin{itemize}
\item \emph{A failed solve is never returned as a solution.} A route that cycles, reaches its iteration limit or meets a singular system returns a result with \code{ok = False}, a \code{reason} code, a one-sentence \code{message} and a \code{suggestion} naming the next route to try. Returning the last iterate with a warning would let a wrong displacement field flow into a design calculation.
\item \emph{Fallback is opt-in.} With \code{fallback=True} the solver tries the suggested routes itself, in the order \texttt{reduced}, \texttt{saddle}, \texttt{condensed}, and lists the routes that failed in \code{fallbacks}; the dense route is only tried when its memory estimate is at most \code{max_dense_bytes}. The default is not to fall back, so that a route that does not suit a problem is noticed.
\item \emph{Warm start.} The result carries \code{active_mask}, and \code{active0=} accepts it, so a sweep over loads or obstacles starts each solve from the previous contact set.
\end{itemize}

\paragraph{A complete example.}
\Cref{lst:ct-static} builds the model problem and solves it. The mesh, the element, the boundary conditions and the assembly are the ones of Chapters~\ref{ch:mesh} and~\ref{ch:elem}; only the last three lines are contact.

\begin{lstlisting}[caption={The membrane obstacle problem with \texttt{Quad4Laplace} and the sparse contact solver.},label={lst:ct-static}]
import numpy as np, scipy.sparse as sp
from fea_engine import FESystem, Quad4Laplace, solve_contact_static_sparse
from fea_engine import mesh as meshlib

m = meshlib.rectangle_mesh(1.0, 1.0, 30, 30)
fes = FESystem(m, Quad4Laplace(), sparse=True)
fes.assemble_stiffness(1.0); fes.assemble_lumped_mass(1.0)
X = m.nodes
edge = np.flatnonzero((X[:, 0] < 1e-9) | (X[:, 0] > 1 - 1e-9) |
                      (X[:, 1] < 1e-9) | (X[:, 1] > 1 - 1e-9))
fes.fix_dofs(edge, ["u"])
free = np.asarray(fes.free_dofs)
K = sp.csr_matrix(fes.K)[free][:, free]
M = sp.diags(sp.csr_matrix(fes.M_lumped).diagonal()[free]).tocsr()
f = -10.0 * (M @ np.ones(len(free)))                       # uniform pressure
c = np.where(np.hypot(X[free, 0] - .5, X[free, 1] - .5) < .2, -0.05, -10.0)
B = sp.identity(len(free), format="csr")                   # every node may touch

r = solve_contact_static_sparse(K, f, B, c)                # route="auto" -> "reduced"
print(r.ok, r.route, r.iterations, len(r.active))
print(r.gap.min(), r.lam.min())                            # both >= 0
\end{lstlisting}

\begin{lstlisting}[numbers=none,language={},keywordstyle={},commentstyle={},frame=single,aboveskip=2pt]
True reduced 2 111
0.0 0.0
\end{lstlisting}

\Cref{lst:ct-fail} shows what a failure looks like and how the fallback chain behaves, by limiting the iteration count so that the \texttt{reduced} route cannot finish.

\begin{lstlisting}[caption={A route that cannot finish reports why and what to try; fallback is opt-in.},label={lst:ct-fail}]
bad = solve_contact_static_sparse(K, f, B, c, route="reduced", max_iter=1)
print(bad.ok, bad.reason, bad.suggestion)
print(bad.message)

r = solve_contact_static_sparse(K, f, B, c, route="reduced", max_iter=1, fallback=True)
print(r.ok, r.route, r.fallbacks)
\end{lstlisting}

\begin{lstlisting}[numbers=none,language={},keywordstyle={},commentstyle={},frame=single,aboveskip=2pt]
False max_iter saddle
route 'reduced' failed (max_iter): no convergence in 1 iterations. Try route='saddle'.
True condensed ['reduced', 'saddle']
\end{lstlisting}

\Cref{tab:ct-map} maps the chapter to the code.

\begin{table}[ht]
\centering\small
\begin{tabularx}{\linewidth}{@{}l X >{\centering\arraybackslash}p{0.11\linewidth}@{}}
\toprule
Concept & Code & Script calls \\
\midrule
sparse primal-dual active set, three routes & \anchor{fea_engine.contact_solvers.solve_contact_static_sparse} & yes \\
result with route, reason, message, suggestion, \texttt{active\_mask} & \anchor{fea_engine.contact_solvers.ContactResult} & no \\
dense reference solvers (\texttt{active\_set}, \texttt{lemke}, \texttt{pgs}, \texttt{pdas}) & \anchor{fea_engine.contact_solvers.solve_contact_static} & yes \\
contact time stepping, first order and BDF2 & \anchor{fea_engine.contact_solvers.solve_contact_dynamics}, \anchor{fea_engine.contact_solvers.ContactDynamicsInfo} & part \\
scalar membrane element & \anchor{fea_engine.elements.scalar.Quad4Laplace} & yes \\
penalty contact elements (the contrast of \Cref{sec:ct-penalty}) & \anchor{fea_engine.elements.contact.GapContactPenalty} & no \\
\bottomrule
\end{tabularx}
\caption{Where the concepts of this chapter live in the repository (tag \booktag). The last column says whether this chapter's own checks call, by name, the code listed in the row: yes (every name), part, or no. Code reached only through another routine counts as no, so the column is a lower bound on what this book exercises.}\label{tab:ct-map}
\end{table}

\section{Verification of the static solvers}\label{sec:ct-verif-static}

A contact solver can return a plausible picture and still be wrong, so the checks are of three kinds: the contact conditions themselves, agreement with independent solvers, and behaviour as the problem grows. All numbers in this section were measured on \machine{}; times are single runs and no test asserts them.

\paragraph{The contact conditions.}
The first check needs no reference solution. For the solution of \Cref{fig:ct-obstacle} the three inequalities and the complementarity product are evaluated directly from the returned $\bm u$ and $\bm\lambda$ (\Cref{sec:ct-problem}); all hold to rounding. The solver applies the same test before it reports success, with a tolerance scaled by the largest displacement and the largest multiplier, so that a converged active set that violates the conditions is reported as a failure with \code{reason="kkt"}.

\paragraph{Agreement with the dense solvers.}
\Cref{tab:ct-agree} compares the two sparse routes, Lemke's method and projected Gauss--Seidel with the dense active-set solver on six different obstacles ($\ctcAgreeN$ unknowns; the number of active constraints ranges from $\ctcAgreeActMin$ to $\ctcAgreeActMax$). The sparse routes and Lemke's method agree with the dense active-set solution to \ctcAgreeRedMax, \ctcAgreeSadMax and \ctcAgreeLemkeMax at most. Projected Gauss--Seidel is an iterative method with its own tolerance and agrees to \ctcAgreePgsMax; the difference is the accuracy of that method, not a disagreement about the active set.

\begin{table}[ht]
\centering\small
\begin{tabular}{lrrrrrrr}
\toprule
obstacle & active & it.\ red. & it.\ sad. & reduced & saddle & Lemke & PGS \\
\midrule
bump & 110 & 7 & 7 & $8.3\times10^{-16}$ & $8.1\times10^{-16}$ & $6.7\times10^{-16}$ & $3.6\times10^{-10}$ \\
plateau $-0.5$ & 60 & 5 & 5 & $8.2\times10^{-16}$ & $7.2\times10^{-16}$ & $2.0\times10^{-16}$ & $2.6\times10^{-10}$ \\
plateau $-0.3$ & 140 & 5 & 5 & $6.8\times10^{-16}$ & $6.8\times10^{-16}$ & $4.3\times10^{-16}$ & $3.7\times10^{-10}$ \\
ridge & 124 & 7 & 7 & $7.2\times10^{-16}$ & $7.3\times10^{-16}$ & $6.7\times10^{-16}$ & $3.5\times10^{-10}$ \\
rough & 20 & 6 & 6 & $6.1\times10^{-16}$ & $5.8\times10^{-16}$ & $2.1\times10^{-16}$ & $4.8\times10^{-11}$ \\
tilted & 72 & 7 & 7 & $8.9\times10^{-16}$ & $7.0\times10^{-16}$ & $3.3\times10^{-16}$ & $3.2\times10^{-10}$ \\
\bottomrule
\end{tabular}

\caption{Six obstacles on the membrane model ($\ctcAgreeN$ unknowns). ``active'' is the number of active constraints; ``it.'' the iterations of each sparse route; the last four columns are the largest difference in $\bm u$ and $\bm\lambda$ to the dense active-set solution.}\label{tab:ct-agree}
\end{table}

\paragraph{Behaviour with problem size.}
\Cref{tab:ct-scaling} and \Cref{fig:ct-time} give the cost of the three routes as the mesh is refined. Two features matter. The iteration count grows slowly with size (from $7$ at the smallest mesh to $\ctcScaleMaxIt$ at $\ctcScaleMaxN$ unknowns), so the cost is dominated by a handful of sparse factorisations. And the dense route grows much faster: at $\ctcScaleDenseN$ unknowns it takes $\ctcScaleDenseT$~s against $\ctcScaleDenseRed$~s (\texttt{reduced}) and $\ctcScaleDenseSad$~s (\texttt{saddle}), and at $\ctcScaleMaxN$ unknowns the sparse routes take $\ctcScaleMaxRed$~s and $\ctcScaleMaxSad$~s. Where the dense route ran, all three routes agree to \ctcScaleDiff.

\begin{table}[ht]
\centering\small
\begin{tabular}{rrrrrrr}
\toprule
$N$ & active & iterations & reduced (s) & saddle (s) & dense (s) & difference to dense \\
\midrule
361 & 84 & 7 & 0.008 & 0.013 & 0.02 & $5.4\times10^{-16}$ \\
841 & 172 & 10 & 0.020 & 0.032 & 0.13 & $9.8\times10^{-16}$ \\
1521 & 295 & 12 & 0.046 & 0.068 & 0.50 & $1.1\times10^{-15}$ \\
2401 & 443 & 13 & 0.078 & 0.109 & 1.27 & $1.2\times10^{-15}$ \\
3481 & 618 & 15 & 0.141 & 0.188 & 3.69 & $1.3\times10^{-15}$ \\
6241 & 1082 & 20 & 0.352 & 0.438 & not run & -- \\
9801 & 1662 & 23 & 0.654 & 0.985 & not run & -- \\
\bottomrule
\end{tabular}

\caption{The model problem on refined meshes. The dense route was not run above 3,500 unknowns. Single runs on \machine.}\label{tab:ct-scaling}
\end{table}

\begin{figure}[ht]
\centering
\includegraphics[width=0.62\linewidth]{figures/ch07_routes_time.pdf}
\caption{Wall time of the three routes against the number of unknowns (data of \Cref{tab:ct-scaling}).}\label{fig:ct-time}
\end{figure}

\paragraph{A larger case, measured on the author's machine.}
The repository script \texttt{fea\_engine/tests/sparse\_contact\_timing.py} was run on the author's machine (12 CPUs) up to a mesh with $\ctcAuthN$ unknowns and $\ctcAuthActive$ active constraints. The \texttt{reduced} route needed $\ctcAuthIt$ iterations and $\ctcAuthRed$~s, the \texttt{saddle} route $\ctcAuthSad$~s, and the two agreed to \ctcAuthDiff. The dense route was not run: by the estimate of \Cref{sec:ct-math} it needs about $\ctcAuthDenseGB$~GB. These numbers were entered by hand from the printed output and are to be regenerated at the pinned release.

\section{Contact dynamics}\label{sec:ct-dyn}

\paragraph{The semi-discrete problem.}
With a mass matrix $\bm M$, the dynamic problem is $\bm M\ddot{\bm u}+\bm K\bm u=\bm f(t)+\bm B^T\bm\lambda$ with the conditions of \eqref{eq:ct-kkt} at every time. A time integrator turns each step into a static contact problem of the form \eqref{eq:ct-kkt} with a modified matrix, so the solvers of \Cref{sec:ct-pdas} are reused unchanged.

\begin{sidenote}{BDF2}
The second-order backward-difference formula advances a step using the two previous ones, $\dot{\bm u}_{n+1}\approx(3\bm u_{n+1}-4\bm u_n+\bm u_{n-1})/(2h)$. It is implicit and stable for stiff problems, and it is second-order accurate on smooth motion. It does not conserve energy at an impact, which is why the energy behaviour in the verification is reported.
\end{sidenote}

\paragraph{The default scheme (first order).}
\fea{} keeps, as the default, the first-order two-step scheme
\begin{equation}\label{eq:ct-twostep}
\bm M(\bm u_{n+1}-2\bm u_n+\bm u_{n-1})+h^2\bm K\bm u_{n+1}=h^2(\bm f_{n+1}+\bm B^T\bm\lambda_{n+1}),\qquad \bm B\bm u_{n+1}\ge\bm c ,
\end{equation}
whose LCP matrix $h^2\bm B(\bm M+h^2\bm K)^{-1}\bm B^T$ is constant and is formed once. Its first step $\bm u_1=\bm u_0+h\bm v_0$ is not checked against the contact conditions. The scheme is first order in $h$.

\paragraph{BDF2.}
The second-order backward differentiation formula~\cite{hairer1996ode2} applied to $\dot{\bm u}=\bm v$ and $\dot{\bm v}=\bm a$ gives $\bm v_{n+1}=(3\bm u_{n+1}-4\bm u_n+\bm u_{n-1})/(2h)$ and $\bm a_{n+1}=(3\bm v_{n+1}-4\bm v_n+\bm v_{n-1})/(2h)$. Substituting into $\bm M\bm a+\bm K\bm u=\bm f+\bm B^T\bm\lambda$ yields, at every step after the first,
\begin{equation}\label{eq:ct-bdf2}
\Bigl(\tfrac{9}{4h^2}\bm M+\bm K\Bigr)\bm u_{n+1}=\bm f_{n+1}+\bm B^T\bm\lambda_{n+1}+\bm M\Bigl[\tfrac{3(4\bm u_n-\bm u_{n-1})}{4h^2}+\tfrac{4\bm v_n-\bm v_{n-1}}{2h}\Bigr],
\end{equation}
together with the contact conditions at $\bm u_{n+1}$. The first step is a BDF1 step, $(\bm M/h^2+\bm K)\bm u_1=\dots$ with $\bm v_1=(\bm u_1-\bm u_0)/h$, and contact is enforced at every step including the first. Both effective matrices are factored once. The unknowns of each step are solved by the static solvers; with \code{method="pdas"} the sparse solver is used at every step with the previous contact set as the warm start.

\begin{algbox}{One BDF2 step with contact}
\begin{pseudo}
\Require $\bm M,\bm K,\bm B,\bm c$, step $h$, states $(\bm u_n,\bm v_n)$ and $(\bm u_{n-1},\bm v_{n-1})$, load $\bm f_{n+1}$, previous active set
\State \textbf{first step (BDF1):} matrix $\bm M/h^2+\bm K$, right-hand side from $\bm u_0,\bm v_0$; \textbf{later steps:} $\tfrac{9}{4h^2}\bm M+\bm K$ and $\bm M\big[\tfrac{3(4\bm u_n-\bm u_{n-1})}{4h^2}+\tfrac{4\bm v_n-\bm v_{n-1}}{2h}\big]$ (factored once)
\State solve $(\ \cdot\ )\bm u_{n+1}=\bm f_{n+1}+\bm B^T\bm\lambda_{n+1}+(\ \cdot\ )$ with $\bm B\bm u_{n+1}\ge\bm c$,\ $\bm\lambda_{n+1}\ge\bm0$, complementarity \Comment{static contact solver, previous contact set as warm start}
\State $\bm v_{n+1}\gets(3\bm u_{n+1}-4\bm u_n+\bm u_{n-1})/(2h)$ \Comment{BDF1: $(\bm u_1-\bm u_0)/h$}
\Ensure $(\bm u_{n+1},\bm v_{n+1},\bm\lambda_{n+1})$ or the failing step with its reason
\end{pseudo}
\end{algbox}

A BDF2 step is one back-substituted contact solve, with both effective matrices factored once. If a contact solve fails inside a run, \code{info} holds \code{failed_step}, \code{reason} and \code{suggestion}, and the energy can rise across contact steps because BDF2 does not conserve it. The code is \anchor{fea_engine.contact_solvers.solve_contact_dynamics} and \anchor{fea_engine.contact_solvers.ContactDynamicsInfo}; \Cref{tab:ct-order} (second order) and \Cref{tab:ct-impact} (impact energy) are the checks.

\begin{lstlisting}[caption={Contact time stepping: default scheme and BDF2.},label={lst:ct-dyn}]
from fea_engine import solve_contact_dynamics
U, Lam = solve_contact_dynamics(M, K, lambda t: f, B, c, 0.005, 200)          # default two-step scheme
U, Lam, info = solve_contact_dynamics(M, K, lambda t: f, B, c, 0.005, 200,
                                      scheme="bdf2", return_info=True)
print(info)     # ContactDynamicsInfo(scheme='bdf2', ok=True, failed_step=None, reason=None)
info.V          # velocities (n_steps + 1, n); None for the two-step scheme
U, Lam = solve_contact_dynamics(M, K, lambda t: f, B, c, 0.005, 200,
                                scheme="bdf2", method="pdas")                 # sparse solve per step
\end{lstlisting}

If a contact solve fails inside a BDF2 run, \code{info} holds the step (\code{failed_step}), the \code{reason} and the \code{suggestion} of the failing solve, and the rows of the output from that step on are \texttt{NaN}, so a failed run cannot be mistaken for a finished one.

\paragraph{Order of accuracy.}
With the obstacle far away the contact is inactive, and the problem reduces to linear dynamics with a known modal solution. \Cref{tab:ct-order} compares the error at $t=2$ of both schemes with the exact modal solution of a free membrane ($\ctcOrdN$ unknowns, lumped mass). The default scheme has observed order $\ctcOrdTwoLast$, that is, order one; BDF2 has observed order $\ctcOrdBdfFirst$ and $\ctcOrdBdfLast$ on the two refinements. The errors differ by a factor $\ctcOrdRatioCoarse$, $\ctcOrdRatioMid$ and $\ctcOrdRatioFine$ at $\Delta t=0.01$, $0.005$ and $0.0025$, so at $\Delta t=0.005$ BDF2 is about $\ctcOrdRatioMid$ times more accurate.

\begin{table}[ht]
\centering\small
\begin{tabular}{rrrrrr}
\toprule
$\Delta t$ & error, two-step & observed order & error, BDF2 & observed order & ratio \\
\midrule
0.01 & $1.6\times10^{-1}$ & -- & $1.2\times10^{-2}$ & -- & 14 \\
0.005 & $8.6\times10^{-2}$ & 0.92 & $3.0\times10^{-3}$ & 1.96 & 29 \\
0.0025 & $4.5\times10^{-2}$ & 0.92 & $7.2\times10^{-4}$ & 2.03 & 63 \\
\bottomrule
\end{tabular}

\caption{Error at $t=2$ of a free membrane against the exact modal solution; contact inactive. ``ratio'' is the error of the two-step scheme divided by that of BDF2.}\label{tab:ct-order}
\end{table}

\begin{figure}[ht]
\centering
\includegraphics[width=0.58\linewidth]{figures/ch07_bdf2_order.pdf}
\caption{Error against time step for the two schemes (data of \Cref{tab:ct-order}).}\label{fig:ct-order}
\end{figure}

\paragraph{Impact: where BDF2 is not enough.}
Order of accuracy is a statement about smooth solutions. An impact is not smooth, so the second test is a bar of 40 nodes that hits a rigid wall at unit speed ($u_0\ge-0.02$), with the energy $\tfrac12\bm v^T\bm M\bm v+\tfrac12\bm u^T\bm K\bm u$ recorded at every step. \Cref{tab:ct-impact} and \Cref{fig:ct-impact} give the result for four time steps. The default scheme never gains energy: its largest energy is $\ctcImpTwoMaxFine$ of the initial one at $\Delta t=0.002$, and it ends at $\ctcImpTwoEndFine$, so it loses more energy than BDF2, which ends at $\ctcImpBdfEndFine$. BDF2 keeps more energy, but its energy rises to $\ctcImpBdfMaxFine$ of the initial value during the impact at $\Delta t=0.002$, and its wall force alternates between large spikes and zero over the first few contacts, where the default scheme gives a smooth force. Neither scheme lets the bar penetrate the wall: the largest penetration in \Cref{tab:ct-impact} is zero for both.

\begin{table}[ht]
\centering\small
\begin{tabular}{rrrrr}
\toprule
$\Delta t$ & scheme & largest energy / initial & final energy / initial & largest penetration \\
\midrule
0.004 & two-step & 1.0000 & 0.9126 & 0 \\
 & BDF2 & 1.0030 & 0.9748 & 0 \\
0.002 & two-step & 1.0000 & 0.9378 & 0 \\
 & BDF2 & 1.0060 & 0.9759 & 0 \\
\bottomrule
\end{tabular}

\caption{Bar impact on a rigid wall: energy relative to the initial energy, and the largest penetration, for the two schemes.}\label{tab:ct-impact}
\end{table}

\begin{figure}[ht]
\centering
\includegraphics[width=0.96\linewidth]{figures/ch07_impact.pdf}
\caption{Bar impact, $\Delta t=0.002$. Left: wall force over the first $0.4$ time units. Right: energy relative to the initial energy over the whole run.}\label{fig:ct-impact}
\end{figure}

\paragraph{How the energy gain depends on the time step.}
The largest BDF2 energy relative to the initial value is \ctcImpBdfMaxVcoarse, \ctcImpBdfMaxCoarse, \ctcImpBdfMaxFine and \ctcImpBdfMaxVfine for $\Delta t=0.008$, $0.004$, $0.002$ and $0.001$ (\Cref{tab:ct-impact}). The trend is monotone, and it runs against the usual expectation: the energy gain \emph{grows} as the step shrinks. The increments fall (the last halving adds only $0.0011$), so the gain tends to a limit near $0.7\%$ and does not keep growing. \Cref{tab:ct-impact-force} gives two further measurements. The same bar without an obstacle, started with a broadband random velocity, keeps \ctcFreeBdfVcoarse, \ctcFreeBdfCoarse, \ctcFreeBdfFine and \ctcFreeBdfVfine of its energy under BDF2 at the four steps, against \ctcFreeTwoVcoarse, \ctcFreeTwoCoarse, \ctcFreeTwoFine and \ctcFreeTwoVfine under the default scheme. The numerical dissipation of BDF2 on high modes therefore falls quickly as $\Delta t$ falls. A likely reading of both results is that, at a large step, this dissipation removes the energy that the contact enforcement adds, and that as the step shrinks the dissipation vanishes and the gain becomes visible. The two measurements agree with that reading, but the gain of the enforcement itself was not isolated (for example by the energy change across each contact step), so it is a likely mechanism and not a demonstrated one. The largest wall force also grows as the step shrinks for both schemes, by about a factor $2$ per halving of $\Delta t$ at the finer steps: the impact is resolved in one or two steps, so the force scales like the impulse divided by $\Delta t$; the BDF2 force is larger than the default scheme's at every step size.

\begin{table}[ht]
\centering\small
\begin{tabular}{rrrrr}
\toprule
 & \multicolumn{2}{c}{largest wall force} & \multicolumn{2}{c}{free bar: energy kept} \\
\cmidrule(lr){2-3}\cmidrule(lr){4-5}
$\Delta t$ & two-step & BDF2 & two-step & BDF2 \\
\midrule
0.008 & 1.21 & 1.95 & 0.0768 & 0.4672 \\
0.004 & 3.36 & 7.83 & 0.1497 & 0.7216 \\
0.002 & 6.49 & 19.65 & 0.2358 & 0.9352 \\
0.001 & 12.86 & 36.93 & 0.3208 & 0.9893 \\
\bottomrule
\end{tabular}

\caption{Largest wall force of the bar impact, and energy kept by a bar without obstacle (broadband initial velocity, $t=1.5$), for both schemes.}\label{tab:ct-impact-force}
\end{table}

BDF2 is therefore more accurate on smooth motion and keeps more energy through an impact, and it is \emph{not} contact-stabilised: the energy gain and the force spikes are real properties of the scheme. Algorithms designed to conserve energy through impact exist~\cite{laursen1997energy}; they are not implemented here. A user who needs a smooth wall force or a guaranteed energy bound should use the default scheme, or inspect the energy history of a BDF2 run, which is what \code{return_info=True} makes possible.

\section{Failure modes and limits}\label{sec:ct-limits}

\begin{limitsbox}
\begin{itemize}
\item \textbf{No convergence guarantee in general.} Finite termination of the active-set iteration is proved for M-matrices; for the stiffness matrix of an elastic solid the method may cycle or stall, and the solver reports it (\texttt{cycled}, \texttt{max\_iter}) instead of returning the last iterate.
\item \textbf{Singular systems.} A \texttt{saddle} system with dependent active rows of $\bm B$, or a $\bm K_{FF}$ that is singular (a body with an unconstrained rigid-body mode), is reported as \texttt{singular}.
\item \textbf{Dense route memory.} The \texttt{condensed} route needs about $8(nm+m^2)$ bytes; it is skipped under \code{fallback=True} when that exceeds \code{max_dense_bytes}, and it was not run at $\ctcAuthN$ unknowns.
\item \textbf{BDF2 is not contact-stabilised.} Energy rises and forces spike during impact (\Cref{tab:ct-impact}); the rise grows as the step shrinks, up to about $0.7\%$ in the steps tested, and no energy result for BDF2 with contact is known to the author.
\item \textbf{Not covered.} Friction in these solvers, Newmark contact dynamics, impact between two flexible bodies, and meshes with $10^5$ or more constraints. Timings are single runs; the large case was run on the author's machine.
\end{itemize}
\end{limitsbox}

\Cref{tab:ct-reasons} lists the reasons the solver can return and what to do about each.

\begin{table}[ht]
\centering\small
\begin{tabularx}{\linewidth}{@{}l X X@{}}
\toprule
\code{reason} & meaning & what to do \\
\midrule
\texttt{not\_applicable} & the \texttt{reduced} route was asked for a $\bm B$ that does not select nodes & use \texttt{saddle} \\
\texttt{cycled} & an active set repeated & try the next route; check that $\bm K$ is positive definite and the obstacle is feasible \\
\texttt{max\_iter} & iteration limit reached & raise \code{max_iter}, warm-start from the returned \code{active_mask}, or try the next route \\
\texttt{singular} & a singular linear system & remove dependent constraints; check the supports of the body \\
\texttt{kkt} & converged active set violates the contact conditions (dense route) & try \code{method="lemke"} of \code{solve_contact_static} \\
\texttt{memory} & the dense route would exceed \code{max_dense_bytes} & raise the limit only if the memory exists; otherwise reformulate \\
\bottomrule
\end{tabularx}
\caption{Failure reasons of \code{solve_contact_static_sparse}.}\label{tab:ct-reasons}
\end{table}

\section{Summary}\label{sec:ct-summary}

We began with a membrane on a bumpy floor and a penalty spring that works until it is stiff enough to be accurate. Frictionless contact is the problem \eqref{eq:ct-kkt}, a convex quadratic program with the inequality $\bm B\bm u\ge\bm c$. A penalty replaces the inequality by a parameter and pays for it in conditioning (\Cref{lst:ct-penalty}, \Cref{fig:ct-penalty}); the primal-dual active-set method treats the inequality exactly as a semismooth Newton iteration and needs a few sparse solves, in the \texttt{reduced} route when $\bm B$ selects nodes and in the \texttt{saddle} route otherwise. The dense condensed form is the reference, and the solver never returns a failed iterate as a solution; it says which route to try next.

In time, BDF2 makes the integration second order on smooth motion, but it is not contact-stabilised, and the energy at an impact is reported and not assumed. The reduced models of Chapter~\ref{ch:hyper} solve the same problem \eqref{eq:ct-kkt} in a reduced space, with the same convention and, for the reduced time stepping, the same BDF2 scheme. That is where Part~III picks it up.

\section*{Exercises}\label{sec:ct-exercises}
\begin{enumerate}
\item \emph{(Implementation)} Run \Cref{lst:ct-static} with a floor at height $-10$ (no contact) and with a floor at $-0.01$ (contact almost everywhere). How does the number of iterations change, and why? (Script: \texttt{fea\_engine/examples/gallery/fea\_09\_sparse\_obstacle\_contact.py}.)
\item \emph{(Implementation)} Repeat the experiment of \Cref{fig:ct-penalty} on a mesh twice as fine. Does the penetration at a fixed $k_p/\overline{K_{ii}}$ change? Explain with the scaling of $\overline{K_{ii}}$ with the mesh.
\item \emph{(Implementation)} Replace $\bm B=\bm I$ by $\bm B=\begin{bmatrix}\bm I&\bm 0\end{bmatrix}+\epsilon\bm S$ with a random sparse $\bm S$, so that \texttt{reduced} is no longer applicable. Which reason does \texttt{reduced} report, and does \texttt{saddle} still match the dense solution? (Test: \texttt{fea\_engine/tests/test\_contact\_solvers\_sparse.py}.)
\item \emph{(Derivation)} Derive \eqref{eq:ct-bdf2} from the two BDF2 formulas given above, and check the first-step matrix $\bm M/h^2+\bm K$ the same way.
\item \emph{(Interpretation)} \Cref{tab:ct-impact} lists four time steps. Add $\Delta t=0.0005$ and decide whether the largest BDF2 energy keeps approaching the limit near $0.7\%$. Then test the explanation of the text by recording the energy change across each contact step (the existing test is \texttt{fea\_engine/tests/test\_contact\_dynamics\_bdf2.py}).
\end{enumerate}

\chapterrefs

\part{Reduced-order modelling (\texttt{rom\_engine})}
\chapter{Projection-based reduction: POD, Galerkin and affine sweeps}\label{ch:pod}

\begin{tcolorbox}[colback=gray!5,colframe=gray!50,title={In this chapter}]
A reduced-order model replaces a system with thousands of unknowns by one with a handful, and is worth building only when the same system has to be solved many times. This chapter builds the simplest such model, the one the rest of Part~III refines: collect solutions (snapshots), extract a basis from them by the singular value decomposition (POD), and project the equations onto it (Galerkin). When the stiffness depends on parameters through a sum of fixed matrices, the projection is done once and every new parameter costs a small dense solve. The chapter measures what this buys and what it does not: how fast the singular values fall, how the error of the reduced solution follows the projection error, whether the choice of snapshots and of the inner product matters, and where the offline cost is repaid.\par\smallskip
\textbf{After this chapter the reader can:} (i) build a POD--Galerkin model with an affine parameter split; (ii) judge it on held-out parameters and not on the training set; (iii) compute the break-even number of queries from its offline and online costs.
\end{tcolorbox}

\section{The problem}\label{sec:pod-problem}

Let $\bm K(\bm\mu)\in\mathbb R^{n\times n}$ be a symmetric positive definite stiffness matrix that depends on parameters $\bm\mu$ (a modulus, a thickness, the position of a load), and let $\bm f(\bm\mu)$ be the load. The full-order model (FOM) is the parametric linear system
\begin{equation}\label{eq:pod-fom}
\bm K(\bm\mu)\,\bm u(\bm\mu)=\bm f(\bm\mu).
\end{equation}
Design studies, optimisation and uncertainty analysis need $\bm u(\bm\mu)$ for hundreds or thousands of values of $\bm\mu$. The reduction method of this chapter rests on one assumption: the solutions $\bm u(\bm\mu)$, which live in $\mathbb R^n$, are well approximated by a subspace of small dimension $r\ll n$. Whether that is true is not a property of the method but of the problem, and \Cref{sec:pod-verif} measures it before anything is built.

\paragraph{The model problem.}
All experiments use the two-region cantilever of the package example \path{rom_engine/examples/two_region_beam_rom.py}: a unit-length Euler--Bernoulli beam of $120$ elements with two stiffness regions, so that $\bm K(\bm\mu)=\mu_1\bm K_1+\mu_2\bm K_2$ with $\mu_1,\mu_2\in[0.5,5]$ and $\bm K_1,\bm K_2$ the stiffness of each region at unit modulus. It has \cpdNFree free unknowns. Two choices of the load matter, and the first is a lesson in itself.

With a single tip load the solution of this beam is $\bm u=\bm u_a/\mu_1+\bm u_b/\mu_2$ for two fixed vectors: the beam is statically determinate, so its solution set has dimension exactly two and two modes are exact. We found this in a first version of the experiments, where the singular values dropped from $10^{2}$ to $10^{-8}$ after the second one and every later question had a trivial answer. That is a useful fact about determinate structures, and it makes a poor test. The experiments below therefore add the \emph{position of a unit point load}, $\xi\in[0.1,0.9]$ of the span, as a third parameter, so that $\bm\mu=(\mu_1,\mu_2,\xi)$. The load vector $\bm f(\xi)$ is not affine in $\xi$; it is assembled by the test script and projected on the fly, which is fine for the accuracy experiments but is not an online-efficient treatment (the limits box says so).

\section{Mathematical structure}\label{sec:pod-math}

\paragraph{Snapshots and POD.}
Solve \eqref{eq:pod-fom} for $N$ training parameters $\bm\mu_1,\dots,\bm\mu_N$ and collect the solutions as the columns of the snapshot matrix $\bm S=[\bm u(\bm\mu_1)\ \cdots\ \bm u(\bm\mu_N)]\in\mathbb R^{n\times N}$. Its singular value decomposition $\bm S=\bm U\bm\Sigma\bm W^T$ gives the basis $\bm V=[\bm U_{:,1}\ \cdots\ \bm U_{:,r}]$ of the first $r$ left singular vectors. By the Eckart--Young theorem no matrix of rank $r$ is closer to $\bm S$ in the Frobenius norm than $\bm V\bm V^T\bm S$, and the distance is
\begin{equation}\label{eq:pod-ey}
\|\bm S-\bm V\bm V^T\bm S\|_F^2=\sum_{k>r}\sigma_k^2 .
\end{equation}
The singular values are therefore a computable bound on how well $r$ modes represent the \emph{training} solutions, and $\sum_{k\le r}\sigma_k^2/\sum_k\sigma_k^2$ is the usual ``energy captured''. The method is known in this form as proper orthogonal decomposition, and in the reduced-basis community as the basis step of the reduced-basis method~\cite{berkooz1993pod,hesthaven2016rb}. Equation~\eqref{eq:pod-ey} says nothing about parameters that were not trained; that is measured by held-out tests in \Cref{sec:pod-verif}.

\paragraph{Weighted POD.}
If the state carries physical units that differ between components (displacements and rotations, as here) or the mesh is non-uniform, the Euclidean inner product weights the unknowns arbitrarily. The mass-weighted variant replaces $\bm S$ by $\bm L^T\bm S$ with $\bm M=\bm L\bm L^T$, takes the SVD, and maps the vectors back, so that $\bm V^T\bm M\bm V=\bm I$ and the basis is optimal in the discrete $L^2$ norm. \anchor{rom_engine.pod.PodBasis} does this with \code{fit(S, M=M)}.

\paragraph{Galerkin projection.}
Approximate $\bm u\approx\bm V\bm q$ and require the residual to be orthogonal to the basis, $\bm V^T(\bm K\bm V\bm q-\bm f)=\bm 0$:
\begin{equation}\label{eq:pod-galerkin}
\bm K_r\bm q=\bm f_r,\qquad \bm K_r=\bm V^T\bm K\bm V\in\mathbb R^{r\times r},\quad \bm f_r=\bm V^T\bm f .
\end{equation}
For symmetric positive definite $\bm K$ this is the orthogonal projection in the energy norm $\|\bm v\|_K^2=\bm v^T\bm K\bm v$, so the reduced solution is the best approximation from $\mathrm{span}(\bm V)$ in that norm, $\|\bm u-\bm V\bm q\|_K=\min_{\bm w}\|\bm u-\bm V\bm w\|_K$ (Céa's lemma applied to the discrete problem). The energy-norm optimality does not transfer to the Euclidean norm used in \cref{eq:pod-ey}, so the Galerkin state error and the Euclidean projection error differ by a factor that depends on $\bm K$; \Cref{sec:pod-verif} measures it. The same projection of a generalised eigenproblem $\bm K\bm\phi=\omega^2\bm M\bm\phi$ gives $\bm K_r\bm z=\omega^2\bm M_r\bm z$, which \code{GalerkinROM.solve_modal} solves (this is the dynamics of Chapter~\ref{ch:dyn} in the subspace).

\tikzfigure[0.8\linewidth]{ch08_projection}{Two projections of the same full-order solution onto the same subspace. The Euclidean projection is what the POD error \eqref{eq:pod-ey} bounds; the Galerkin solution minimises the energy-norm error instead, so the two points differ and so do the two errors. Schematic, not computed.}{fig:pod-projection}

\paragraph{Affine parameter dependence.}
Projecting \eqref{eq:pod-galerkin} costs $O(n^2r)$ for a dense matrix, and if it were repeated for every $\bm\mu$ nothing would be saved. When
\begin{equation}\label{eq:pod-affine}
\bm K(\bm\mu)=\sum_{q=1}^{Q}\theta_q(\bm\mu)\,\bm K_q,
\end{equation}
with scalar functions $\theta_q$ and parameter-free matrices $\bm K_q$, the projection commutes with the sum, $\bm K_r(\bm\mu)=\sum_q\theta_q(\bm\mu)\,\bm V^T\bm K_q\bm V$, so each $\bm V^T\bm K_q\bm V$ is computed once offline and the online cost for a new $\bm\mu$ is the sum of $Q$ small matrices and an $r\times r$ solve, independent of $n$. \anchor{rom_engine.affine.AffineDecomposition} stores the $\bm K_q$ and the function $\bm\theta$, \code{project(V)} does the offline step, and \code{assemble_reduced(mu)} the online sum.

\begin{algbox}{Proper orthogonal decomposition of a snapshot matrix (optionally mass-weighted)}
\begin{pseudo}
\Require snapshots $\bm S\in\mathbb R^{n\times N}$, optional mass matrix $\bm M$, rank $r$ or energy tolerance $\epsilon$
\If{$\bm M$ is given}
  \State $\bm L\gets\mathrm{chol}(\bm M)$ with $\bm M=\bm L\bm L^T$;\ $\tilde{\bm S}\gets\bm L^T\bm S$
\Else
  \State $\tilde{\bm S}\gets\bm S$
\EndIf
\State $[\bm U,\bm\Sigma,\bm W]\gets\mathrm{svd}(\tilde{\bm S})$ (economy size)
\State $r\gets$ the given rank, or the smallest $r$ with $\sum_{k\le r}\sigma_k^2/\sum_k\sigma_k^2\ge1-\epsilon$
\State $\bm V\gets\bm U_{:,1:r}$;\ \textbf{if} $\bm M$ is given \textbf{then} $\bm V\gets\bm L^{-T}\bm V$ \Comment{now $\bm V^T\bm M\bm V=\bm I$}
\Ensure basis $\bm V\in\mathbb R^{n\times r}$ and the singular values $\sigma_k$
\end{pseudo}
\begin{algdet}
\item[Cost] SVD of an $n\times N$ matrix, $O(nN^2)$ for $N\le n$; one Cholesky factorisation of $\bm M$ if weighted.
\item[Fails when] the singular values decay slowly (a large $r$ is needed); the snapshots do not span the parameter range of interest; the captured energy is high but the held-out error is not (it says nothing about untrained parameters).
\item[Code] \anchor{rom_engine.pod.PodBasis}, \code{fit(S, M=M)}.
\item[Verified in] \Cref{fig:pod-pod} (singular values and held-out error), \Cref{sec:pod-verif}.
\end{algdet}
\end{algbox}

\begin{algbox}{Offline--online reduced basis for a parametric stiffness}
\begin{pseudo}
\Require affine terms $\{\bm K_q\}_{q=1}^Q$ and coefficient function $\bm\theta(\bm\mu)$, load $\bm f$, training size $N$, rank $r$
\State \textbf{Offline (once):}
\State choose $\bm\mu_1,\dots,\bm\mu_N$ by an optimal Latin hypercube \Comment{\code{optimal_lhs}}
\For{$i=1,\dots,N$}
  \State solve $\bm K(\bm\mu_i)\bm u_i=\bm f$ \Comment{$N$ full-order solves}
\EndFor
\State $\bm V\gets$ POD of $[\bm u_1\ \cdots\ \bm u_N]$ \Comment{previous algorithm}
\For{$q=1,\dots,Q$}
  \State $\bm K_{q,r}\gets\bm V^T\bm K_q\bm V$
\EndFor
\State $\bm f_r\gets\bm V^T\bm f$
\State \textbf{Online (for each new $\bm\mu$):}
\State $\bm K_r\gets\sum_q\theta_q(\bm\mu)\bm K_{q,r}$;\ solve $\bm K_r\bm q=\bm f_r$
\State $\bm u\gets\bm V\bm q$ \Comment{only if the full field is needed}
\Ensure reduced coordinates $\bm q$ (and $\bm u$) for every query
\end{pseudo}
\begin{algdet}
\item[Cost] Offline: $N$ full solves, one SVD, $Q$ projections of $O(n^2r)$ each for dense $\bm K_q$. Online: $O(Qr^2+r^3)$, independent of $n$; the lift adds $O(nr)$.
\item[Fails when] $\bm K$ or $\bm f$ is not a sum of fixed terms (the online step then depends on $n$); the training range does not cover the query range; a loading that moves, as in \Cref{sec:pod-verif}, makes the solution set larger than the basis.
\item[Code] \anchor{rom_engine.affine.AffineDecomposition} (\code{project}, \code{assemble_reduced}), \anchor{rom_engine.galerkin.GalerkinROM}.
\item[Verified in] \Cref{tab:pod-cost} and \Cref{fig:pod-sweep} (costs, break-even), \Cref{sec:pod-verif} (held-out error).
\end{algdet}
\end{algbox}

\tikzfigure{ch08_pipeline}{The offline--online pipeline of \Cref{lst:pod-run}. Blue: the full-order model; violet: the reduced-order model; grey: plain data; green: error measure; dashed: optional or validation only. Only the lift touches the full dimension $n$ online.}{fig:pod-pipeline}

\section{Implementation in the packages}\label{sec:pod-impl}

\Cref{lst:pod-run} is the whole pipeline on the model problem; \Cref{tab:pod-map} lists where each step lives. The three classes are deliberately separate: \code{PodBasis} knows only about snapshots, \code{GalerkinROM} only about a basis and matrices, \code{AffineDecomposition} only about a sum of matrices. A basis can come from any source (POD, modal truncation, the Krylov bases of Chapter~\ref{ch:bt}) and the other two classes do not care.

\begin{lstlisting}[caption={Offline--online pipeline on the two-region beam (matrices K1, K2 are the dense free-dof blocks; theta is the coefficient function).},label={lst:pod-run}]
import numpy as np
from rom_engine import PodBasis, GalerkinROM, AffineDecomposition
from rom_engine.sampling import optimal_lhs

# offline
rng = np.random.default_rng(0)
mus = 0.5 + 4.5 * optimal_lhs(16, 2, rng=rng)          # (EI1, EI2) training design
S = np.column_stack([np.linalg.solve(m[0] * K1 + m[1] * K2, F) for m in mus])
pod = PodBasis().fit(S, n_modes=8)                       # or M=M for mass-weighted POD
aff = AffineDecomposition([K1, K2], lambda mu: [mu[0], mu[1]]).project(pod.V)
Fr = pod.V.T @ F

# online: one 8 x 8 solve per parameter
q = np.linalg.solve(aff.assemble_reduced((2.0, 1.0)), Fr)
u = pod.V @ q

# modal analysis in the subspace
rom = GalerkinROM(pod.V).reduce_system(2.0 * K1 + 1.0 * K2, M)
freq_hz, _, _ = rom.solve_modal(n_modes=5)
\end{lstlisting}

\begin{table}[ht]
\centering\small
\begin{tabularx}{\linewidth}{@{}l X >{\centering\arraybackslash}p{0.11\linewidth}@{}}
\toprule
Concept & Code & Script calls \\
\midrule
snapshot SVD, truncation by rank or energy, mass-weighted variant & \anchor{rom_engine.pod.PodBasis} & yes \\
Galerkin projection of $\bm K,\bm M,\bm f$; static and modal solves & \anchor{rom_engine.galerkin.GalerkinROM} & yes \\
affine sum $\sum_q\theta_q\bm K_q$, offline projection, online assembly & \anchor{rom_engine.affine.AffineDecomposition} & yes \\
optimal Latin hypercube design & \anchor{rom_engine.sampling.optimal_lhs} & yes \\
parametric mode sensitivities (used for mode-tracking in parameter studies) & \anchor{rom_engine.parametric_modes.mode_sensitivities} & no \\
worked examples & \path{rom_engine/examples/gallery/rom_01_pod_basis.py}, \path{rom_02_galerkin_parametric_rom.py} & -- \\
\bottomrule
\end{tabularx}
\caption{Where the concepts of this chapter live in the repository (release \booktag). The last column says whether the chapter's script \texttt{scripts/chNN\_assets.py} calls, by name, the code listed in the row: yes (every name), part, or no. Code reached only through another routine counts as no, so the column is a lower bound on what this book exercises.}\label{tab:pod-map}
\end{table}

\section{Verification}\label{sec:pod-verif}

Each claim below is a measurement from \texttt{scripts/ch08\_assets.py}, run on \machine{}. Errors are relative, $\|\bm u-\bm V\bm q\|/\|\bm u\|$ in the Euclidean norm, and are the worst value over \cpdNTest held-out parameters drawn at random, not seen in training. Timings are single runs; no test asserts them.

\subsection{Singular values and the held-out error}

\Cref{fig:pod-pod} uses \cpdNSnap training solutions from an optimal Latin hypercube. The singular values fall by a factor of \cpdSvDecade between the first and the fourth, and reach \cpdSvTwelve of the first at $r=12$. The held-out error follows them down more slowly: Galerkin state error \cpdGalerkinErrSix with $6$ modes, \cpdGalerkinErrTwelve with $12$, and \cpdGalerkinErrLast with $24$. To reach $5\%$ takes \cpdModesForFivePct modes, and to reach $1\%$ takes \cpdModesForOnePct. The training spectrum goes to rounding level (\cpdSvLast of the first value by the last of the $60$), but the held-out error does not, which is the point of testing on parameters that were not trained.

\begin{figure}[ht]
\centering
\includegraphics[width=0.98\linewidth]{figures/ch08_pod.pdf}
\caption{Left: normalised singular values of the \cpdNSnap training snapshots. Right: worst relative error on \cpdNTest unseen parameters against the number of modes $r$, for the Euclidean projection error and for the Galerkin solve with Euclidean and with mass-weighted POD. Produced by \texttt{scripts/ch08\_assets.py}.}\label{fig:pod-pod}
\end{figure}

The Galerkin state error is never more than \cpdGalerkinOverProj times the Euclidean projection error of the same basis (largest ratio over the $24$ sizes), consistent with the energy-norm optimality of \cref{eq:pod-galerkin}: the reduced solve loses little against the best approximation in the subspace. A singular-value plot alone would have suggested a much faster decay than the unseen parameters show, and the measured \cpdModesForOnePct modes for $1\%$ is the figure to trust.

\subsection{Weighted against Euclidean POD}

Mass-weighted POD has $\|\bm V^T\bm M\bm V-\bm I\|_{\max}$ equal to \cpdMassOrth{} and Euclidean POD has $\|\bm V^T\bm V-\bm I\|_{\max}$ equal to \cpdEucOrth{}, so each basis is orthonormal in its own inner product to rounding. On this problem the two choices give nearly the same Galerkin error (\cpdMassErrSix against \cpdGalerkinErrSix at $r=6$, and \cpdMassErrTwelve against \cpdGalerkinErrTwelve at $r=12$); the weighted basis is lower at \cpdMassWins of the \cpdMassCompared sizes. We read this as a small effect on a beam with a uniform mesh, and we do not claim it carries over to non-uniform meshes or to meshes with mixed unknown types, where the weighting is expected to matter more; that was not tested.

\subsection{How the snapshots are chosen}

\Cref{fig:pod-snap} compares an optimal (maximin) Latin hypercube with uniform random sampling at equal budget, with the basis size set to $\min(N,12)$ and five random repetitions for each size. The optimal design has the lower median error at \cpdOlhsBetterCount of the \cpdSnapCases sizes: at $N=20$ the median is \cpdSnapOlhsB against \cpdSnapRandB for random sampling, and at $N=10$ it is \cpdSnapOlhsA against \cpdSnapRandA. Five repetitions are few, and the one size where random wins shows that the ordering is not guaranteed; the dependence on the budget is the stronger result, since the error falls with $N$ until the basis size saturates at $12$.

\begin{figure}[ht]
\centering
\includegraphics[width=0.62\linewidth]{figures/ch08_snapshots.pdf}
\caption{Median Galerkin state error on \cpdNTest unseen parameters against the number of snapshots, optimal Latin hypercube against random sampling, five repetitions each.}\label{fig:pod-snap}
\end{figure}

\subsection{Modal analysis in the subspace}

The basis above was built from \emph{static} solutions, and a static basis is not a modal basis. \Cref{tab:pod-modal} projects $\bm K$ and $\bm M$ at $\bm\mu=(2,1)$ and compares the lowest frequencies with the full eigenproblem. The first frequency is within \cpdModalErrOneAtSix with $6$ modes and the third within \cpdModalErrThreeAtSix, but the fifth is still off by \cpdModalErrFiveAtTwelve with $12$ modes. A basis built for the displacement under these loads resolves the modes that the loads excite, and the higher modes require either a richer set of snapshots or a different method (Chapters~\ref{ch:bt} and~\ref{ch:cms}).

\begin{table}[ht]
\centering\small
\begin{tabular}{rrrrr}
\toprule
modes $r$ & mode 1 & mode 2 & mode 3 & mode 5 \\
\midrule
4 & $2.7\times10^{-4}$ & $1.9\times10^{-2}$ & $1.9\times10^{-1}$ & -- \\
6 & $2.7\times10^{-4}$ & $4.2\times10^{-3}$ & $1.8\times10^{-2}$ & $3.1\times10^{-1}$ \\
8 & $1.3\times10^{-4}$ & $3.0\times10^{-3}$ & $1.2\times10^{-2}$ & $5.3\times10^{-2}$ \\
12 & $6.1\times10^{-5}$ & $1.4\times10^{-3}$ & $7.5\times10^{-3}$ & $3.9\times10^{-2}$ \\
\bottomrule
\end{tabular}

\caption{Relative error of the lowest frequencies of the reduced eigenproblem against the full model, for bases of $r$ POD modes of the static snapshots (a dash: fewer than five modes in the basis).}\label{tab:pod-modal}
\end{table}

\subsection{Offline and online cost}

\Cref{fig:pod-sweep} and \Cref{tab:pod-cost} compare one parameter query in three ways on two sizes of the same beam: re-assembling the full matrix with the helper of the example and solving, summing the stored affine terms and solving, and the reduced solve of \cref{lst:pod-run}. The reduced solve takes about \cpdOnlineUs~$\mu$s per query in the largest run, independent of $n$ in the two sizes measured, while a full query takes \cpdCostSpeedAffine times longer at \cpdNFree unknowns with affine assembly and \cpdCostSpeedReasm times longer with re-assembly, and \cpdCostSpeedAffineBig and \cpdCostSpeedReasmBig times longer at \cpdBigNf unknowns. The offline step (snapshots, SVD, projection) is repaid after $3$ queries at the smaller size and \cpdBreakEvenBig at the larger one against re-assembly, and after about \cpdBreakEvenAffBig queries at the larger size against the affine full solve.

\begin{figure}[ht]
\centering
\includegraphics[width=0.62\linewidth]{figures/ch08_sweep.pdf}
\caption{Time per query: reduced solve, full solve with affine assembly, full solve with re-assembly, at two sizes of the beam. Logarithmic axis; single runs on \machine.}\label{fig:pod-sweep}
\end{figure}

\begin{table}[ht]
\centering\small
\begin{tabular}{rrrrrrr}
\toprule
 & & offline & \multicolumn{3}{c}{one query ($\mu$s)} & \\
\cmidrule(lr){4-6}
elements & dofs & (ms) & reduced & full, affine & full, re-assembled & break-even queries \\
\midrule
120 & 240 & 11 & 17 & 656 & 4305 & 3 \\
600 & 1200 & 742 & 18 & 46508 & 88376 & 8 \\
\bottomrule
\end{tabular}

\caption{Offline cost and per-query cost on two sizes of the beam; break-even is the offline cost divided by the saving per query against re-assembly. Single runs on \machine.}\label{tab:pod-cost}
\end{table}

\paragraph{What the figures include.} The online time above is the reduced solve only. Lifting the reduced answer to the full space, $\bm u=\bm V\bm q$, is a separate cost of \cpdRecUs~$\mu$s at \cpdNFree{} unknowns and \cpdRecUsBig~$\mu$s at \cpdBigNf; it is paid whenever the full field is needed and is not paid if only a few reduced outputs are. The offline time of \Cref{tab:pod-cost} includes everything that precedes the first query (snapshots, SVD, projection), and the break-even counts include it. A workflow figure that includes it, for $\cpdWorkflowQ$ queries with lifting, is a factor \cpdWorkflow{} against full re-assembly at \cpdNFree{} unknowns and \cpdWorkflowBig{} at \cpdBigNf, and \cpdWorkflowSparseFree{} and \cpdWorkflowSparseFreeBig{} against the dense solve with affine assembly: smaller than the per-query ratios, as it must be. The error is reported twice. Over the test parameters the worst relative error of the full state is \cpdGlobErr{} (\cpdNFree{} unknowns) and \cpdGlobErrBig{} (\cpdBigNf), and the worst relative error of the tip deflection, the local quantity of interest, is \cpdLocErr{} and \cpdLocErrBig. They coincide here because the tip load makes the tip the dominant degree of freedom of the response, so this test does not show whether a local error can exceed the global one; the three-parameter study of \Cref{sec:pod-verif}, with a load that moves, is the better place to look for that, and a local quantity such as a reaction or a peak stress was not computed.

\tikzfigure[0.7\linewidth]{ch08_breakeven}{Cumulative cost of $Q$ queries. The reduced model pays $T_{\mathrm{off}}$ first and less per query afterwards; the lines cross at the break-even $Q^\ast$. Schematic: measured values are in \Cref{tab:pod-cost}.}{fig:pod-breakeven}

These speed-ups are inflated by the baseline. The full solves here use dense \code{numpy.linalg.solve} on a matrix whose entries are almost all zero, which is how the example builds its models; a banded or sparse factorisation of this tri-diagonal-block system is far cheaper than the dense solve, and the speed-up against it would be much smaller. The result that survives is the break-even structure (a modest offline cost repaid after tens of queries) and the scaling (the reduced solve does not grow with $n$), not the ratios.

\section{Limits and failure modes}\label{sec:pod-limits}

\begin{limitsbox}
\begin{itemize}
\item \textbf{The basis is only as good as the snapshots.} The held-out error (\Cref{fig:pod-pod}) is the only guide to accuracy at untrained parameters; singular values describe training data. No certified error bound is given in this chapter (Chapter~\ref{ch:freq} adds one).
\item \textbf{Affine structure is a requirement, not a convenience.} If $\bm K$ or $\bm f$ is not a sum of fixed terms, the online step is not independent of $n$. Here the load position makes $\bm f$ non-affine, and the experiments of \Cref{sec:pod-verif} project it per query, so they test accuracy and not online efficiency; the cost experiments use the fixed tip load.
\item \textbf{Geometry changes are not handled.} A parameter that moves the mesh breaks \cref{eq:pod-affine}; the package has no empirical-interpolation route for it in this chapter.
\item \textbf{Static bases do not give modal accuracy} (\Cref{tab:pod-modal}); the higher frequencies are poorly captured.
\item \textbf{Determinate structures are trivial.} A single-load statically determinate beam has an exactly two-dimensional solution set; conclusions drawn from it do not transfer.
\item \textbf{Weighted POD} was compared on one uniform-mesh beam only.
\item \textbf{Timings} are single runs on \machine{} with a dense baseline; sizes up to \cpdBigNf unknowns only.
\end{itemize}
\end{limitsbox}

\section{Summary}\label{sec:pod-summary}

POD extracts a low-dimensional basis from solutions by the SVD, \cref{eq:pod-ey}; Galerkin projection turns the large system into an $r\times r$ one, \cref{eq:pod-galerkin}, and is the best approximation in the energy norm. If the parameter dependence is affine, \cref{eq:pod-affine}, the projection is done once and each new parameter costs a small solve. On the model beam, held-out error is the figure to trust: \cpdModesForOnePct modes for $1\%$ although the singular values suggest fewer. The offline cost is repaid after a few to a few tens of queries, but the speed-up ratios reported here depend on a dense baseline.

\section*{Exercises}\label{sec:pod-exercises}
\begin{enumerate}
\item \emph{(Implementation)} Reproduce \Cref{fig:pod-pod} with only $(\mu_1,\mu_2)$ varying and the tip load fixed. Where do the singular values drop to rounding level, and why? (Script: \texttt{scripts/ch08\_assets.py}.)
\item \emph{(Implementation)} Replace the held-out parameter set by one drawn from $\xi\in[0.1,0.3]$ only. Does a basis trained on the full range still reach the same error?
\item \emph{(Implementation)} Use \code{fit(S, energy_threshold=0.9999)} instead of a fixed rank. How many modes does it select, and what is the held-out error? Does the training energy predict the held-out error?
\item \emph{(Implementation)} Add a third affine term (a stiffer segment) to \cref{lst:pod-run}. What changes in the offline projection cost and in the online cost?
\item \emph{(Implementation)} Use a sparse factorisation (\code{scipy.sparse.linalg.splu}) for the full solves of \Cref{tab:pod-cost} and report the new break-even numbers.
\item \emph{(Derivation)} For snapshots $\bm S$ with the mass inner product, show that the squared projection error of the snapshot set on the leading $r$ POD modes equals the sum of the discarded squared singular values, $\sum_{j>r}\sigma_j^2$.
\item \emph{(Interpretation)} In this chapter the held-out error follows the projection error. Describe a situation in which the Galerkin model would be much worse than its projection error, and what that would signal.
\end{enumerate}

\chapterrefs

\chapter{Systems-and-control reduction}\label{ch:bt}
\stubchapter{Chapter 9}

\chapter{Frequency-domain ROMs and certified error bounds}\label{ch:freq}
\stubchapter{Chapter 10}

\chapter{Substructuring and component reduction}\label{ch:cms}
\stubchapter{Chapter 11}

\chapter{Non-intrusive identification and nonlinear ROMs}\label{ch:nonint}
\stubchapter{Chapter 12}

\chapter{Hyper-reduction and contact reduced-order models}\label{ch:hyper}
\stubchapter{Chapter 13}

\part{Verification}
\chapter{Verification, reproduction studies and a parametric pipeline}\label{ch:verif}

\begin{tcolorbox}[colback=gray!5,colframe=gray!50,title={In this chapter}]
Chapters~\ref{ch:arch} to~\ref{ch:hyper} each ended with their own verification section. This chapter collects what those sections established, says what kind of evidence each one is and what it cannot show, then goes further in two directions: the reconstruction of published studies from the packages alone, and one complete run, from mesh to a held-out prediction, that is reported according to Appendix~\ref{ch:conv}.
\textbf{After this chapter the reader can:} (i) place any number of this book on the ladder from unit test to integrated run and say what it does not prove; (ii) read the outcome table of the reproduction studies, including the ones that failed; (iii) run the integrated pipeline on their own model and report it with offline, online and lifting cost and with a local quantity next to the global error.
\end{tcolorbox}

\section{What ``verified'' means in this book}\label{sec:vf-levels}

A claim that a piece of code is correct can be supported at six levels, and each level answers a different question. \Cref{fig:vf-pyramid} lists them. A unit test (level 1) asks whether the code does what its author meant on one input. A closed-form benchmark (level 2) asks whether that is also what the mathematics says, on a case where the exact answer is known and is not computed by the code under test. Convergence orders and the agreement of independent solvers (level 3) ask whether the error falls at the rate the method promises, and whether two different routes to the same quantity give the same number. For reduced models (level 4) the question changes to accuracy on parameters or loads that were not used in training, and to whether a bound held. A reproduction (level 5) asks whether a method described in somebody else's words, rebuilt from the packages without their code, recovers what its authors report. An integrated run (level 6) asks whether the pieces still work when they are connected: a model built by the mesh tools, solved by the engine, reduced, and used on data that none of the earlier steps saw.

\tikzfigure[0.98\linewidth]{ch14_pyramid}{The six levels of evidence used in this book. Lower levels contain many small cases and are cheap to repeat; higher levels contain few cases and test more of the stack at once. A higher level does not replace a lower one: it can pass while a component below it is wrong in a region it does not visit.}{fig:vf-pyramid}

Two rules apply to every level. The first is that the reference must not come from the code under test (Appendix~\ref{ch:conv}); a closed form, a quadrature of an integral, a much finer model computed another way, or a published number are references, the same solver at another setting is not. The second is that a pass is a statement about the cases that were run. The sections below therefore state, for each group of results, the case, the reference, the number, and what was not tried.

\section{Component verification in one table}\label{sec:vf-component}

\Cref{tab:vf-map} lists, chapter by chapter, one headline check of the verification sections. The numbers are the macros that the chapter scripts generated (they are the same numbers as in the chapters, not new measurements) on the machine named in each chapter. The last column points to the section where the case, the reference and the qualifications are given; the entry in the table is a summary and should not be quoted without them.

\begin{table}[tbp]
\centering\footnotesize
\setlength{\tabcolsep}{4pt}
\caption{Headline checks from Chapters~\ref{ch:elem} to~\ref{ch:hyper}. ``Level'' refers to \Cref{fig:vf-pyramid}. An error without a unit is a relative error.}\label{tab:vf-map}
\begin{tabularx}{\linewidth}{@{}l c >{\raggedright\arraybackslash}p{3.6cm} >{\raggedright\arraybackslash}X >{\raggedright\arraybackslash}p{2.0cm}@{}}
\toprule
Ch. & Level & Check and reference & Result & Where \\
\midrule
\ref{ch:elem} & 2 & patch test, all \cecPatchCount{} continuum elements, exact linear field & worst error \cecPatchWorst{} & \Cref{sec:el-verif} \\
\ref{ch:elem} & 2, 3 & Euler--Bernoulli cantilever frequency, closed form; Timoshenko tip deflection, closed form & observed orders \cecEBFreqOrder{} (theory 4) and \cecTimoOrder{} (theory 2) & \Cref{sec:el-verif} \\
\ref{ch:dyn} & 3 & Newmark, error against the exact discrete solution & observed order \cdyNewmarkOrder{} (theory 2) & \Cref{ch:dyn} \\
\ref{ch:dyn} & 3 & harmonic response against an independent modal sum & difference \cdyHarmDiff{} & \Cref{ch:dyn} \\
\ref{ch:nonlin} & 2 & tangent stiffness against finite differences of the internal force & tangent error \cnlTangentWorst{} & \Cref{sec:nl-verif} \\
\ref{ch:nonlin} & 2 & von Mises truss, closed form (displacement control; arc length) & errors \cnlTrussErrDC{} and \cnlTrussErrAL{}; load control fails at the limit load & \Cref{sec:nl-verif} \\
\ref{ch:batch} & 3 & chunked batched assembly against the unchunked batched assembly & difference \cbkChunkDiff{} & \Cref{sec:bk-verif} \\
\ref{ch:contact} & 3 & contact problems on six obstacles (\ctcAgreeN{} unknowns), sparse routes, Lemke and projected Gauss--Seidel against the dense active-set solver & differences at most \ctcAgreeRedMax{}, \ctcAgreeSadMax{} and \ctcAgreeLemkeMax{}; projected Gauss--Seidel \ctcAgreePgsMax{} (its own tolerance) & \Cref{sec:ct-verif-static} \\
\ref{ch:contact} & 3 & BDF2 contact, observed order & \ctcOrdBdfLast{} (theory 2); largest energy reaches \ctcImpBdfMaxVfine{} times the initial value at the finest step & \Cref{ch:contact} \\
\ref{ch:pod} & 4 & POD--Galerkin, held-out state error & \cpdGalerkinErrTwelve{} with 12 modes, \cpdGalerkinErrLast{} with 24 & \Cref{sec:pod-verif} \\
\ref{ch:bt} & 4 & a priori balanced-truncation bound against the grid maximum of the error & held in \cbtBoundHolds{} of \cbtBoundCases{} cases; the miss is at rounding level (\cbtMissErr{} against \cbtMissBound{}) & \Cref{sec:bt-verif} \\
\ref{ch:freq} & 4 & certified bound, linear-programme lower bound & held at all \cfrLpBoundFinite{} frequencies where it was finite; median bound-to-error ratio \cfrLpRatioMed{} & \Cref{sec:fr-verif} \\
\ref{ch:cms} & 4 & Craig--Bampton, largest error of the first four frequencies & \ccmCbFiveErr{} with five modes, \ccmCbAllErr{} with all (exact) & \Cref{sec:cm-verif} \\
\ref{ch:nonint} & 3, 4 & Duffing backbone by harmonic balance against time integration & frequency error \cnmErrOne{}, \cnmErrThree{}, \cnmErrSeven{} with 1, 3, 7 harmonics & \Cref{sec:ni-verif} \\
\ref{ch:hyper} & 4 & reduced contact problem with the SVD dual basis, against the full-order BDF2 reference & reduced KKT residual \chdSvdTwentyKktOut{} (20 modes), yet negative lifted multipliers; the NMF basis gives none: a small residual is not a valid contact state & \Cref{sec:hy-verif} \\
\bottomrule
\end{tabularx}
\end{table}

Three entries deserve a second look because they show the limits of the table itself. The balanced-truncation bound is reported as held in \cbtBoundHolds{} of \cbtBoundCases{} cases, not in all of them: the one miss is a case at rounding level, where both numbers are far below anything that matters, and the book reports it as a miss. The certified bound of Chapter~\ref{ch:freq} is valid wherever it is finite, and it is finite at few frequencies, so its guarantee covers a small part of the band. And the SVD dual basis of Chapter~\ref{ch:hyper} satisfies the reduced KKT conditions to rounding error and still gives negative lifted multipliers, which is why that chapter judges a reduced contact model by the lifted multipliers and not by the reduced residual.

\subsection{What the component checks do not show}\label{sec:vf-notshow}

The checks of \Cref{tab:vf-map} cover the element types and the analyses named, on simple geometries: beams, a plate, small plane bodies, one cantilever with two material regions. They do not cover three-dimensional plasticity against an independent reference, shells in contact, or the Koiter--Newton solvers; those parts are exercised by the unit tests of the repository and are not evaluated here. The reduced models are tested on one or two structures each, with parameters and loads drawn from the ranges of the training. Error is reported against a full-order model computed with the same engine, so a defect common to both (a wrong element, a wrong load) would not be seen by a ROM check; the closed-form checks of levels 2 and 3 are what guard against that. Wall-clock times are for the 2-CPU machine named in each chapter and for a dense baseline unless a chapter says otherwise.

\stubchapter{Sections 14.3 to 14.5: reproduction studies, integrated parametric pipeline, limits and exercises}


\appendix
\chapter{Differentiable assembly and autograd tangents}\label{ch:diff}
\stubchapter{Appendix A}


\backmatter
\end{document}
